%============================================================
%  Bd. 11 - Ordnungsrelationen
%============================================================

\documentclass[main.tex]{subfiles}

\ifSubfilesClassLoaded{
  \BandSetupStandalone{B11}
}{
}

\title{Bd. 11 - Ordnungsrelationen}
\author{Martin Kunze}
\date{}
\setcounter{file}{11}

\begin{document}

\title{Bd. 11 - Ordnungsrelationen}
\author{Martin Kunze}
\date{}
\hypersetup{pageanchor=false}
\maketitle
\hypersetup{pageanchor=true}
\setcounter{tocdepth}{3}
\tableofcontents
\setcounter{file}{11}
\renewcommand{\FormulaBandID}{11}

\chapter{Einleitung}

Dieser Band behandelt partielle, totale und wohlgeordnete Relationsoperatoren
unabhängig von einem besonderen Zahlbereich. Minimum, Maximum, Schranken,
Infimum und Supremum werden hier als allgemeine Ordnungsbegriffe entwickelt.
Natürliche Ordnungen und halbverbandsinduzierte Ordnungen erscheinen
anschließend als Anwendungen in den jeweiligen Fachbänden.

\chapter{Ordnungsrelationen}

\section{Begriff der Ordnung}

\subsection{Axiome einer partiellen Ordnung}

\FormulaAxiomDeltaK[Antisymmetrie]{%
  x\leq y \dsep y\leq x \vdash x=y%
}{OrderAntisymmetry}{
  \DeltaRow{Mengen}{A \dsep x \dsep y}
  \DeltaRow{Relationsoperatoren}{\leq}
}

\subsection{Definition der partiellen Ordnung}

\FormulaDefDeltaK[Begriff der partiellen Ordnung]{\PartOrd{A,\leq}}{Partielle Ordnung}{
  \DeltaRow{Mengen}{A \dsep x \dsep y \dsep z}
  \DeltaRow{Relationsoperatoren}{\leq}
  \DeltaRow{\textbf{Axiome}}{}
  \DeltaRow{Binärer Relationsoperator}
           {\leq\ \text{auf }A}
  \DeltaRow{Reflexivität}
           {x\in A \vdash x\leq x}
           [\FormulaRefAuto{EqRelReflexivity}[ax]]
  \DeltaRow{Transitivität}
           {x\leq y \dsep y\leq z \vdash x\leq z}
           [\FormulaRefAuto{EqRelTransitivity}[ax]]
  \DeltaRow{Antisymmetrie}
           {x\leq y \dsep y\leq x \vdash x=y}
           [\FormulaRefAuto{x\leq y \dsep y\leq x \vdash x=y}[ax]]
  \DeltaRow{\textbf{Neue Symbole}}{}
  \DeltaPrem{Partielle Ordnung}{\PartOrd{A,\leq}}
}

Die Ordnung wird somit stets zusammen mit ihrem Träger angegeben. Ihre
Auswertung erfolgt in Infixnotation: \(x\leq y\), nicht durch die
Zugehörigkeit eines geordneten Paares zu einer Relationsmenge.

\subsection{Dualer Relationsoperator einer partiellen Ordnung}

\FormulaDefDeltaK[Dualer Relationsoperator auf einem Träger]{%
  x,y\in A\vdash
  \bigl(x\geq y\leftrightarrow y\leq x\bigr)
}{OrderDualRelation}{%
  \DeltaRow{Mengen}{A\dsep x\dsep y}
  \DeltaRow{Relationsoperatoren}{\leq\dsep\geq}
  \DeltaRow{Auswertung}{x,y\in A\vdash(x\geq y\leftrightarrow y\leq x)}
  \DeltaRow{\textbf{Neues Symbol}}{}
  \DeltaPrem{Dualer Relationsoperator}{\geq}
}

\FormulaThmDeltaK[Der duale Relationsoperator ist eine partielle Ordnung]{%
  \PartOrd{A,\leq}\vdash\PartOrd{A,\geq}
}{OrderDualPartialOrder}{%
  \DeltaRow{Mengen}{A\dsep x\dsep y\dsep z}
  \DeltaRow{Relationsoperatoren}{\leq\dsep\geq}
}
\begin{tabproofsplitwide}
  \proofpartwideindR[Reflexivität]{%
    \PartOrd{A,\leq}\dsep x\in A\vdash x\geq x
  }{%
    \PartOrd{A,\leq}\dsep x\in A\vdash x\geq x
  }[OrderDualReflexive]
    \proofstepwidestar[1]{\PartOrd{A,\leq}}{\rA}
    \proofstepwidestar[2]{x\in A}{\rA}
    \proofstepwidestar[1,2]{x\leq x}{%
      \FormulaRefAuto{EqRelReflexivity}[ax]{2}}
    \proofstepwidestar[2]{x\geq x\leftrightarrow x\leq x}{%
      \FormulaRefAuto{OrderDualRelation}[def]{2,2}}
    \proofstepwidestar[1,2]{x\geq x}{%
      \FormulaRefAuto{P\leftrightarrow Q, Q\vdash P}{4,3}}
  \closeproofpartwideindR

  \proofpartwideindR[Transitivität]{%
    \begin{aligned}[t]
      &\PartOrd{A,\leq}\dsep x,y,z\in A
        \dsep x\geq y\dsep y\geq z\vdash{}\\
      &x\geq z
    \end{aligned}
  }{%
    \begin{aligned}[t]
      &\PartOrd{A,\leq}\dsep x,y,z\in A
        \dsep x\geq y\dsep y\geq z\vdash{}\\
      &x\geq z
    \end{aligned}
  }[OrderDualTransitive]
    \proofstepwidestar[1]{\PartOrd{A,\leq}}{\rA}
    \proofstepwidestar[2]{x\in A}{\rA}
    \proofstepwidestar[3]{y\in A}{\rA}
    \proofstepwidestar[4]{z\in A}{\rA}
    \proofstepwidestar[5]{x\geq y}{\rA}
    \proofstepwidestar[6]{y\geq z}{\rA}
    \proofstepwidestar[2,3,5]{y\leq x}{%
      \FormulaRefAuto{P\leftrightarrow Q, P\vdash Q}{%
        \FormulaRefAuto{OrderDualRelation}[def]{2,3},5}}
    \proofstepwidestar[3,4,6]{z\leq y}{%
      \FormulaRefAuto{P\leftrightarrow Q, P\vdash Q}{%
        \FormulaRefAuto{OrderDualRelation}[def]{3,4},6}}
    \proofstepwidestar[1,2,3,4,5,6]{z\leq x}{%
      \FormulaRefAuto{EqRelTransitivity}[ax]{8,7}}
    \proofstepwidestar[2,4]{x\geq z\leftrightarrow z\leq x}{%
      \FormulaRefAuto{OrderDualRelation}[def]{2,4}}
    \proofstepwidestar[1,2,3,4,5,6]{x\geq z}{%
      \FormulaRefAuto{P\leftrightarrow Q, Q\vdash P}{10,9}}
  \closeproofpartwideindR

  \proofpartwideindR[Antisymmetrie]{%
    \begin{aligned}[t]
      &\PartOrd{A,\leq}\dsep x,y\in A
        \dsep x\geq y\dsep y\geq x\vdash x=y
    \end{aligned}
  }{%
    \begin{aligned}[t]
      &\PartOrd{A,\leq}\dsep x,y\in A
        \dsep x\geq y\dsep y\geq x\vdash x=y
    \end{aligned}
  }[OrderDualAntisymmetric]
    \proofstepwidestar[1]{\PartOrd{A,\leq}}{\rA}
    \proofstepwidestar[2]{x\in A}{\rA}
    \proofstepwidestar[3]{y\in A}{\rA}
    \proofstepwidestar[4]{x\geq y}{\rA}
    \proofstepwidestar[5]{y\geq x}{\rA}
    \proofstepwidestar[2,3,4]{y\leq x}{%
      \FormulaRefAuto{P\leftrightarrow Q, P\vdash Q}{%
        \FormulaRefAuto{OrderDualRelation}[def]{2,3},4}}
    \proofstepwidestar[2,3,5]{x\leq y}{%
      \FormulaRefAuto{P\leftrightarrow Q, P\vdash Q}{%
        \FormulaRefAuto{OrderDualRelation}[def]{3,2},5}}
    \proofstepwidestar[1,2,3,4,5]{x=y}{%
      \FormulaRefAuto{x\leq y \dsep y\leq x\vdash x=y}[ax]{7,6}}
  \closeproofpartwideindR

  \proofpartwide{Schluss}
    \proofstepwidestar[1]{\PartOrd{A,\leq}}{\rA}
    \proofstepwidestar[1]{x\in A\vdash x\geq x}{%
      \FormulaRefAuto{OrderDualReflexive}[thm]{1}}
    \proofstepwidestar[1]{%
      x,y,z\in A\dsep x\geq y\dsep y\geq z\vdash x\geq z}{%
      \FormulaRefAuto{OrderDualTransitive}[thm]{1}}
    \proofstepwidestar[1]{%
      x,y\in A\dsep x\geq y\dsep y\geq x\vdash x=y}{%
      \FormulaRefAuto{OrderDualAntisymmetric}[thm]{1}}
    \proofstepwidestar[1]{\PartOrd{A,\geq}}{%
      \FormulaRefAuto{Partielle Ordnung}[def]{2,3,4}}
  \closeproofpartwide
\end{tabproofsplitwide}

\subsection{Axiom einer totalen Ordnung}

\FormulaAxiomDelta[Vergleichbarkeit]{%
  x\in A \dsep y\in A \vdash x\leq y \lor y\leq x%
}{
  \DeltaRow{Mengen}{A \dsep x \dsep y}
  \DeltaRow{Relationsoperatoren}{\leq}
}

\subsection{Definition der totalen Ordnung}

\FormulaDefDeltaK[Begriff der totalen Ordnung]{\TotOrd{A,\leq}}{Totale Ordnung}{
  \DeltaRow{Mengen}{A \dsep x \dsep y \dsep z}
  \DeltaRow{Relationsoperatoren}{\leq}
  \DeltaRow{\textbf{Axiome}}{}
  \DeltaRow{Partielle Ordnung}
           {\PartOrd{A,\leq}}
           [\FormulaRefAuto{Partielle Ordnung}]
  \DeltaRow{Vergleichbarkeit}
           {\begin{aligned}[t]
              &x\in A\dsep y\in A\vdash{}\\[-2pt]
              &x\leq y \lor y\leq x
            \end{aligned}}
           [\FormulaRefAuto{x\in A \dsep y\in A
             \vdash x\leq y \lor y\leq x}[ax]]
  \DeltaRow{\textbf{Neue Symbole}}{}
  \DeltaPrem{Totale Ordnung}{\TotOrd{A,\leq}}
}

% ============================================================
% Kleinste / groesste Elemente, Minimum / Maximum
% (auf Basis einer partiellen Ordnung)
% ============================================================

% ============================================================
% Minimum / Maximum (als eindeutige Terme via \iota)
% Kleinstes/Groesstes Element als Spezialfall T=A
% ============================================================

\subsection{Extremale Elemente}

\subsubsection{Minimales Element und Minimum}

Ein minimales Element besitzt in der Teilmenge kein echt kleineres Element.
Ein Minimum ist demgegenüber ein kleinstes Element und liegt unter
\emph{jedem} Element der Teilmenge. Die folgenden Sätze halten den genauen
Zusammenhang fest.

\FormulaDefDeltaK[Minimales Element einer Teilmenge]{%
  T\subseteq A\vdash
  \operatorname{MinEl}_{A,\leq}(m,T)
  \coloneqq
  m\in T\land\forall x\in T\,(x\leq m\rightarrow x=m)
}{OrderMinimalElement}{%
  \DeltaRow{Mengen}{A\dsep T\dsep m\dsep x}
  \DeltaRow{Relationsoperatoren}{\leq}
  \DeltaPrem{Partielle Ordnung}{\PartOrd{A,\leq}}
  \DeltaRow{Teilmenge}{T\subseteq A}
  \DeltaRow{\textbf{Neues Symbol}}{}
  \DeltaPrem{Minimales Element}{\operatorname{MinEl}_{A,\leq}(m,T)}
}

% ------------------------------------------------------------
% (1) Eindeutigkeit des Minimum-Kriteriums
% ------------------------------------------------------------
\FormulaThmDeltaK[Eindeutigkeit eines Minimums]{%
  \exists m\in T\,\forall x\in T\,(m\leq x)
  \vdash \exists! m\in T\,\forall x\in T\,(m\leq x)%
}{MinimumUnique}{
  \DeltaRow{Mengen}{A \dsep T \dsep m \dsep n \dsep x}
  \DeltaRow{Relationsoperatoren}{\leq}
  \DeltaPrem{Partielle Ordnung}{\PartOrd{A,\leq}}
}
\begin{tabproof}
  \proofstep{1}{\exists m\in T\,\forall x\in T\,(m\leq x)}{\rA}

  \proofstep{2}{m\in T \land \forall x\in T\,(m\leq x)}{\rA}
  \proofstep{3}{n\in T \land \forall x\in T\,(n\leq x)}{\rA}

  \proofstep{2}{m\in T}{\rAEa{2}}
  \proofstep{2}{\forall x\in T\,(m\leq x)}{\rAEb{2}}
  \proofstep{3}{n\in T}{\rAEa{3}}
  \proofstep{3}{\forall x\in T\,(n\leq x)}{\rAEb{3}}

  \proofstep{2,3}{m\leq n}{\rRE{6,\rUE{5}}}
  \proofstep{2,3}{n\leq m}{\rRE{4,\rUE{7}}}

  \proofstep{2,3}{m=n}{%
    \FormulaRefAuto{x\leq y \dsep y\leq x\vdash x=y}[ax]{8,9}}

  \proofstep{1}{\exists! m\in T\,\forall x\in T\,(m\leq x)}{%
    \UEI{1,2,3,10}}
\end{tabproof}


% ------------------------------------------------------------
% (2) Definition: Minimum als \iota-Term (partielle Definition)
% ------------------------------------------------------------
\FormulaDefDeltaK[Minimum]{%
  \begin{aligned}
    &\exists m\in T\,\forall x\in T\,(m\leq x)\\
    &\vdash \min_{A,\leq}(T)\coloneqq
      \iota m\in T\,\forall x\in T\,(m\leq x)
  \end{aligned}%
}{Minimum}{
  \DeltaRow{Mengen}{A \dsep T \dsep m \dsep x}
  \DeltaRow{Relationsoperatoren}{\leq}
  \DeltaPrem{Partielle Ordnung}{\PartOrd{A,\leq}}
}

\FormulaThmDeltaK[Ein Minimum ist ein minimales Element]{%
  \begin{aligned}[t]
    &T\subseteq A\dsep
      \exists q\in T\,\forall x\in T\,(q\leq x)\vdash{}\\
    &\operatorname{MinEl}_{A,\leq}\bigl(\min_{A,\leq}(T),T\bigr)
  \end{aligned}
}{MinimumIsMinimalElement}{%
  \DeltaRow{Mengen}{A\dsep T\dsep q\dsep x}
  \DeltaRow{Relationsoperatoren}{\leq}
  \DeltaPrem{Partielle Ordnung}{\PartOrd{A,\leq}}
}
\begin{tabproofwide}
  \proofstepwidestar[1]{T\subseteq A}{\rA}
  \proofstepwidestar[2]{%
    \exists q\in T\,\forall x\in T\,(q\leq x)}{\rA}
  \proofstepwidestar[2]{\min_{A,\leq}(T)\in T}{%
    \FormulaRefAuto{Minimum}[def]{2}}
  \proofstepwidestar[4]{x\in T}{\rA}
  \proofstepwidestar[5]{x\leq\min_{A,\leq}(T)}{\rA}
  \proofstepwidestar[2,4]{\min_{A,\leq}(T)\leq x}{%
    \FormulaRefAuto{Minimum}[def]{2,4}}
  \proofstepwidestar[2,4,5]{x=\min_{A,\leq}(T)}{%
    \FormulaRefAuto{x\leq y\dsep y\leq x\vdash x=y}[ax]{5,6}}
  \proofstepwidestar[2,4]{%
    x\leq\min_{A,\leq}(T)\rightarrow x=\min_{A,\leq}(T)}{%
    \rRI{5,7}}
  \proofstepwidestar[2]{%
    \forall x\in T\,
      \bigl(x\leq\min_{A,\leq}(T)\rightarrow x=\min_{A,\leq}(T)\bigr)}{%
    \rUI{\rRI{4,8}}}
  \proofstepwidestar[1,2]{%
    \operatorname{MinEl}_{A,\leq}\bigl(\min_{A,\leq}(T),T\bigr)}{%
    \FormulaRefAuto{OrderMinimalElement}[def]{1,\rAI{3,9}}}
\end{tabproofwide}

\subsubsection{Maximales Element und Maximum}

\FormulaDefDeltaK[Maximales Element einer Teilmenge]{%
  T\subseteq A\vdash
  \operatorname{MaxEl}_{A,\leq}(m,T)
  \coloneqq
  m\in T\land\forall x\in T\,(m\leq x\rightarrow x=m)
}{OrderMaximalElement}{%
  \DeltaRow{Mengen}{A\dsep T\dsep m\dsep x}
  \DeltaRow{Relationsoperatoren}{\leq}
  \DeltaPrem{Partielle Ordnung}{\PartOrd{A,\leq}}
  \DeltaRow{Teilmenge}{T\subseteq A}
  \DeltaRow{\textbf{Neues Symbol}}{}
  \DeltaPrem{Maximales Element}{\operatorname{MaxEl}_{A,\leq}(m,T)}
}

% ------------------------------------------------------------
% (3) Eindeutigkeit des Maximum-Kriteriums (kurz mit  \exists m\in T ...)
% ------------------------------------------------------------
\FormulaThmDeltaK[Eindeutigkeit eines Maximums]{%
  \exists m\in T\,\forall x\in T\,(x\leq m)
  \vdash \exists! m\in T\,\forall x\in T\,(x\leq m)%
}{MaximumUnique}{
  \DeltaRow{Mengen}{A \dsep T \dsep m \dsep n \dsep x}
  \DeltaRow{Relationsoperatoren}{\leq}
  \DeltaPrem{Partielle Ordnung}{\PartOrd{A,\leq}}
}
\begin{tabproof}
  \proofstep{1}{\exists m\in T\,\forall x\in T\,(x\leq m)}{\rA}

  \proofstep{2}{m\in T \land \forall x\in T\,(x\leq m)}{\rA}
  \proofstep{3}{n\in T \land \forall x\in T\,(x\leq n)}{\rA}

  \proofstep{2}{m\in T}{\rAEa{2}}
  \proofstep{2}{\forall x\in T\,(x\leq m)}{\rAEb{2}}
  \proofstep{3}{n\in T}{\rAEa{3}}
  \proofstep{3}{\forall x\in T\,(x\leq n)}{\rAEb{3}}

  \proofstep{2,3}{m\leq n}{\rRE{4,\rUE{7}}}
  \proofstep{2,3}{n\leq m}{\rRE{6,\rUE{5}}}

  \proofstep{2,3}{m=n}{%
    \FormulaRefAuto{x\leq y \dsep y\leq x\vdash x=y}[ax]{8,9}}

  \proofstep{1}{\exists! m\in T\,\forall x\in T\,(x\leq m)}{%
    \UEI{1,2,3,10}}
\end{tabproof}

% ------------------------------------------------------------
% (4) Definition: Maximum als \iota-Term (partielle Definition)
% ------------------------------------------------------------
\FormulaDefDeltaK[Maximum]{%
  \begin{aligned}
    &\exists m\in T\,\forall x\in T\,(x\leq m)\\
    &\vdash \max_{A,\leq}(T)\coloneqq
      \iota m\in T\,\forall x\in T\,(x\leq m)
  \end{aligned}%
}{Maximum}{
  \DeltaRow{Mengen}{A \dsep T \dsep m \dsep x}
  \DeltaRow{Relationsoperatoren}{\leq}
  \DeltaPrem{Partielle Ordnung}{\PartOrd{A,\leq}}
}

\FormulaThmDeltaK[Ein Maximum ist ein maximales Element]{%
  \begin{aligned}[t]
    &T\subseteq A\dsep
      \exists q\in T\,\forall x\in T\,(x\leq q)\vdash{}\\
    &\operatorname{MaxEl}_{A,\leq}\bigl(\max_{A,\leq}(T),T\bigr)
  \end{aligned}
}{MaximumIsMaximalElement}{%
  \DeltaRow{Mengen}{A\dsep T\dsep q\dsep x}
  \DeltaRow{Relationsoperatoren}{\leq}
  \DeltaPrem{Partielle Ordnung}{\PartOrd{A,\leq}}
}
\begin{tabproofwide}
  \proofstepwidestar[1]{T\subseteq A}{\rA}
  \proofstepwidestar[2]{%
    \exists q\in T\,\forall x\in T\,(x\leq q)}{\rA}
  \proofstepwidestar[2]{\max_{A,\leq}(T)\in T}{%
    \FormulaRefAuto{Maximum}[def]{2}}
  \proofstepwidestar[4]{x\in T}{\rA}
  \proofstepwidestar[5]{\max_{A,\leq}(T)\leq x}{\rA}
  \proofstepwidestar[2,4]{x\leq\max_{A,\leq}(T)}{%
    \FormulaRefAuto{Maximum}[def]{2,4}}
  \proofstepwidestar[2,4,5]{x=\max_{A,\leq}(T)}{%
    \FormulaRefAuto{x\leq y\dsep y\leq x\vdash x=y}[ax]{6,5}}
  \proofstepwidestar[2,4]{%
    \max_{A,\leq}(T)\leq x\rightarrow x=\max_{A,\leq}(T)}{%
    \rRI{5,7}}
  \proofstepwidestar[2]{%
    \forall x\in T\,
      \bigl(\max_{A,\leq}(T)\leq x\rightarrow x=\max_{A,\leq}(T)\bigr)}{%
    \rUI{\rRI{4,8}}}
  \proofstepwidestar[1,2]{%
    \operatorname{MaxEl}_{A,\leq}\bigl(\max_{A,\leq}(T),T\bigr)}{%
    \FormulaRefAuto{OrderMaximalElement}[def]{1,\rAI{3,9}}}
\end{tabproofwide}

\FormulaThmDeltaK[In einer totalen Ordnung ist jedes minimale Element ein Minimum]{%
  \begin{aligned}[t]
    &\TotOrd{A,\leq}\dsep T\subseteq A\dsep
      \operatorname{MinEl}_{A,\leq}(m,T)\vdash{}\\
    &m=\min_{A,\leq}(T)
  \end{aligned}
}{TotalOrderMinimalElementIsMinimum}{%
  \DeltaRow{Mengen}{A\dsep T\dsep m\dsep x}
  \DeltaRow{Relationsoperatoren}{\leq}
}
\begin{tabproofwide}
  \proofstepwidestar[1]{\TotOrd{A,\leq}}{\rA}
  \proofstepwidestar[2]{T\subseteq A}{\rA}
  \proofstepwidestar[3]{\operatorname{MinEl}_{A,\leq}(m,T)}{\rA}
  \proofstepwidestar[1]{\PartOrd{A,\leq}}{%
    \FormulaRefAuto{Totale Ordnung}[def]{1}}
  \proofstepwidestar[2,3]{%
    \begin{aligned}[t]
      &m\in T\land{}\\[-2pt]
      &\forall x\in T\,(x\leq m\rightarrow x=m)
    \end{aligned}}{%
    \FormulaRefAuto{OrderMinimalElement}[def]{2,3}}
  \proofstepwidestar[2,3]{m\in T}{\rAEa{5}}
  \proofstepwidestar[7]{x\in T}{\rA}
  \proofstepwidestar[2,3]{m\in A}{%
    \FormulaRefAuto{A\subseteq B,\,x\in A\vdash x\in B}{2,6}}
  \proofstepwidestar[2,7]{x\in A}{%
    \FormulaRefAuto{A\subseteq B,\,x\in A\vdash x\in B}{2,7}}
  \proofstepwidestar[1,2,3,7]{m\leq x\lor x\leq m}{%
    \FormulaRefAuto{x\in A\dsep y\in A
      \vdash x\leq y\lor y\leq x}[ax]{8,9}}
  \proofstepwidestar[11]{m\leq x}{\rA}
  \proofstepwidestar[12]{x\leq m}{\rA}
  \proofstepwidestar[2,3,7,12]{x=m}{%
    \rRE{12,\rUE{\rAEb{5}}}}
  \proofstepwidestar[2,3,7,12]{m\leq x}{%
    \rIE{13,\FormulaRefAuto{EqRelReflexivity}[ax]{9}}}
  \proofstepwidestar[1,2,3,7]{m\leq x}{%
    \rOE{10,11,11,12,14}}
  \proofstepwidestar[1,2,3]{\forall x\in T\,(m\leq x)}{%
    \rUI{\rRI{7,15}}}
  \proofstepwidestar[1,2,3]{%
    \begin{aligned}[t]
      &\exists q\in T\,\bigl(\\[-2pt]
      &\qquad\forall x\in T\,(q\leq x)\bigr)
    \end{aligned}}{%
    \rEI{\rAI{6,16}}}
  \proofstepwidestar[1,2,3]{%
    \begin{aligned}[t]
      &\exists! q\in T\,\bigl(\\[-2pt]
      &\qquad\forall x\in T\,(q\leq x)\bigr)
    \end{aligned}}{%
    \FormulaRefAuto{MinimumUnique}[thm]{17}}
  \proofstepwidestar[1,2,3]{%
    \begin{aligned}[t]
      &\min_{A,\leq}(T)\in T\land{}\\[-2pt]
      &\forall x\in T\,(\min_{A,\leq}(T)\leq x)
    \end{aligned}}{%
    \rAI{\FormulaRefAuto{Minimum}[def]{17},
         \FormulaRefAuto{Minimum}[def]{17}}}
  \proofstepwidestar[1,2,3]{m=\min_{A,\leq}(T)}{%
    \FormulaRefAuto{\exists! x\,P(x),\;P(a),\;P(b)\vdash a=b}{%
      18,\rAI{6,16},19}}
\end{tabproofwide}

\FormulaThmDeltaK[In einer totalen Ordnung ist jedes maximale Element ein Maximum]{%
  \begin{aligned}[t]
    &\TotOrd{A,\leq}\dsep T\subseteq A\dsep
      \operatorname{MaxEl}_{A,\leq}(m,T)\vdash{}\\
    &m=\max_{A,\leq}(T)
  \end{aligned}
}{TotalOrderMaximalElementIsMaximum}{%
  \DeltaRow{Mengen}{A\dsep T\dsep m\dsep x}
  \DeltaRow{Relationsoperatoren}{\leq}
}
\begin{tabproofwide}
  \proofstepwidestar[1]{\TotOrd{A,\leq}}{\rA}
  \proofstepwidestar[2]{T\subseteq A}{\rA}
  \proofstepwidestar[3]{\operatorname{MaxEl}_{A,\leq}(m,T)}{\rA}
  \proofstepwidestar[1]{\PartOrd{A,\leq}}{%
    \FormulaRefAuto{Totale Ordnung}[def]{1}}
  \proofstepwidestar[2,3]{%
    \begin{aligned}[t]
      &m\in T\land{}\\[-2pt]
      &\forall x\in T\,(m\leq x\rightarrow x=m)
    \end{aligned}}{%
    \FormulaRefAuto{OrderMaximalElement}[def]{2,3}}
  \proofstepwidestar[2,3]{m\in T}{\rAEa{5}}
  \proofstepwidestar[7]{x\in T}{\rA}
  \proofstepwidestar[2,3]{m\in A}{%
    \FormulaRefAuto{A\subseteq B,\,x\in A\vdash x\in B}{2,6}}
  \proofstepwidestar[2,7]{x\in A}{%
    \FormulaRefAuto{A\subseteq B,\,x\in A\vdash x\in B}{2,7}}
  \proofstepwidestar[1,2,3,7]{x\leq m\lor m\leq x}{%
    \FormulaRefAuto{x\in A\dsep y\in A
      \vdash x\leq y\lor y\leq x}[ax]{9,8}}
  \proofstepwidestar[11]{x\leq m}{\rA}
  \proofstepwidestar[12]{m\leq x}{\rA}
  \proofstepwidestar[2,3,7,12]{x=m}{%
    \rRE{12,\rUE{\rAEb{5}}}}
  \proofstepwidestar[2,3,7,12]{x\leq m}{%
    \rIE{13,\FormulaRefAuto{EqRelReflexivity}[ax]{9}}}
  \proofstepwidestar[1,2,3,7]{x\leq m}{%
    \rOE{10,11,11,12,14}}
  \proofstepwidestar[1,2,3]{\forall x\in T\,(x\leq m)}{%
    \rUI{\rRI{7,15}}}
  \proofstepwidestar[1,2,3]{%
    \begin{aligned}[t]
      &\exists q\in T\,\bigl(\\[-2pt]
      &\qquad\forall x\in T\,(x\leq q)\bigr)
    \end{aligned}}{%
    \rEI{\rAI{6,16}}}
  \proofstepwidestar[1,2,3]{%
    \begin{aligned}[t]
      &\exists! q\in T\,\bigl(\\[-2pt]
      &\qquad\forall x\in T\,(x\leq q)\bigr)
    \end{aligned}}{%
    \FormulaRefAuto{MaximumUnique}[thm]{17}}
  \proofstepwidestar[1,2,3]{%
    \begin{aligned}[t]
      &\max_{A,\leq}(T)\in T\land{}\\[-2pt]
      &\forall x\in T\,(x\leq\max_{A,\leq}(T))
    \end{aligned}}{%
    \rAI{\FormulaRefAuto{Maximum}[def]{17},
         \FormulaRefAuto{Maximum}[def]{17}}}
  \proofstepwidestar[1,2,3]{m=\max_{A,\leq}(T)}{%
    \FormulaRefAuto{\exists! x\,P(x),\;P(a),\;P(b)\vdash a=b}{%
      18,\rAI{6,16},19}}
\end{tabproofwide}

Die Umkehrungen der beiden vorigen allgemeinen Sätze gelten in partiellen
Ordnungen nicht: Sind \(a\) und \(b\) unvergleichbar, dann sind in
\(\{a,b\}\) beide Elemente minimal und maximal, die Teilmenge besitzt aber
weder ein Minimum noch ein Maximum. In totalen Ordnungen beseitigt die
Vergleichbarkeit genau dieses Hindernis.


% ============================================================
% Schranken
% ============================================================

\subsection{Schranken}

\subsubsection{Untere und obere Schrankenmengen}

\FormulaDefDeltaK[Untere Schrankenmenge]{%
  T\subseteq A \vdash
  \LB_{A,\leq}(T)\coloneqq
  \{\, a\in A \mid \forall x\in T\,\bigl(a\leq x\bigr)\,\}%
}{LowerBoundSet}{
  \DeltaRow{Mengen}{A\dsep T \dsep a \dsep x}
  \DeltaRow{Relationsoperatoren}{\leq}
  \DeltaPrem{Partielle Ordnung}{\PartOrd{A,\leq}}
}

\FormulaDefDeltaK[Obere Schrankenmenge]{%
  T\subseteq A \vdash
  \UB_{A,\leq}(T)\coloneqq
  \{\, a\in A \mid \forall x\in T\,\bigl(x\leq a\bigr)\,\}%
}{UpperBoundSet}{
  \DeltaRow{Mengen}{A\dsep T \dsep a \dsep x}
  \DeltaRow{Relationsoperatoren}{\leq}
  \DeltaPrem{Partielle Ordnung}{\PartOrd{A,\leq}}
}

\subsubsection{Infimum und Supremum}

\FormulaDefDeltaK[Infimum]{%
  \begin{aligned}
    &T\subseteq A \dsep \exists m\in \LB_{A,\leq}(T)\,
      \forall a\in \LB_{A,\leq}(T)\,\bigl(a\leq m\bigr)\\
    &\vdash \inf_{A,\leq}(T)\coloneqq \max_{A,\leq}\bigl(\LB_{A,\leq}(T)\bigr)
  \end{aligned}%
}{Infimum}{
  \DeltaRow{Mengen}{A\dsep T \dsep m \dsep a}
  \DeltaRow{Relationsoperatoren}{\leq}
  \DeltaPrem{Partielle Ordnung}{\PartOrd{A,\leq}}
}

\FormulaDefDeltaK[Supremum]{%
  \begin{aligned}
    &T\subseteq A \dsep \exists m\in \UB_{A,\leq}(T)\,
      \forall a\in \UB_{A,\leq}(T)\,\bigl(m\leq a\bigr)\\
    &\vdash \sup_{A,\leq}(T)\coloneqq \min_{A,\leq}\bigl(\UB_{A,\leq}(T)\bigr)
  \end{aligned}%
}{Supremum}{
  \DeltaRow{Mengen}{A\dsep T \dsep m \dsep a}
  \DeltaRow{Relationsoperatoren}{\leq}
  \DeltaPrem{Partielle Ordnung}{\PartOrd{A,\leq}}
}

\FormulaThmDeltaK[Charakterisierung der oberen Schrankenmenge]{%
  T\subseteq A\vdash
  \Bigl(u\in\UB_{A,\leq}(T)\leftrightarrow
    \bigl(u\in A\land\forall x\in T\,(x\leq u)\bigr)\Bigr)
}{UpperBoundSetCharacterization}{%
  \DeltaRow{Mengen}{A\dsep T\dsep u\dsep x}
  \DeltaRow{Relationsoperatoren}{\leq}
  \DeltaPrem{Partielle Ordnung}{\PartOrd{A,\leq}}
}
\begin{tabproofwide}
  \proofstepwidestar[1]{T\subseteq A}{\rA}
  \proofstepwidestar[1]{%
    u\in\UB_{A,\leq}(T)\leftrightarrow
    \bigl(u\in A\land\forall x\in T\,(x\leq u)\bigr)}{%
    \FormulaRefAuto{UpperBoundSet}[def]{1}}
\end{tabproofwide}

\FormulaThmDeltaK[Auswertung der Zugehörigkeit zur oberen Schrankenmenge]{%
  T\subseteq A\dsep u\in\UB_{A,\leq}(T)\vdash
  u\in A\land\forall x\in T\,(x\leq u)
}{UpperBoundSetElim}{%
  \DeltaRow{Mengen}{A\dsep T\dsep u\dsep x}
  \DeltaRow{Relationsoperatoren}{\leq}
  \DeltaPrem{Partielle Ordnung}{\PartOrd{A,\leq}}
}
\begin{tabproofwide}
  \proofstepwidestar[1]{T\subseteq A}{\rA}
  \proofstepwidestar[2]{u\in\UB_{A,\leq}(T)}{\rA}
  \proofstepwidestar[1]{%
    u\in\UB_{A,\leq}(T)\leftrightarrow
      \bigl(u\in A\land\forall x\in T\,(x\leq u)\bigr)}{%
    \FormulaRefAuto{UpperBoundSetCharacterization}[thm]{1}}
  \proofstepwidestar[1,2]{u\in A\land\forall x\in T\,(x\leq u)}{%
    \FormulaRefAuto{P\leftrightarrow Q, P\vdash Q}{3,2}}
\end{tabproofwide}

\FormulaThmDeltaK[Einführung in die obere Schrankenmenge]{%
  T\subseteq A\dsep u\in A\dsep
  \forall x\in T\,(x\leq u)\vdash u\in\UB_{A,\leq}(T)
}{UpperBoundSetIntro}{%
  \DeltaRow{Mengen}{A\dsep T\dsep u\dsep x}
  \DeltaRow{Relationsoperatoren}{\leq}
  \DeltaPrem{Partielle Ordnung}{\PartOrd{A,\leq}}
}
\begin{tabproofwide}
  \proofstepwidestar[1]{T\subseteq A}{\rA}
  \proofstepwidestar[2]{u\in A}{\rA}
  \proofstepwidestar[3]{\forall x\in T\,(x\leq u)}{\rA}
  \proofstepwidestar[2,3]{u\in A\land\forall x\in T\,(x\leq u)}{%
    \rAI{2,3}}
  \proofstepwidestar[1]{%
    u\in\UB_{A,\leq}(T)\leftrightarrow
      \bigl(u\in A\land\forall x\in T\,(x\leq u)\bigr)}{%
    \FormulaRefAuto{UpperBoundSetCharacterization}[thm]{1}}
  \proofstepwidestar[1,2,3]{u\in\UB_{A,\leq}(T)}{%
    \FormulaRefAuto{P\leftrightarrow Q, Q\vdash P}{5,4}}
\end{tabproofwide}

\FormulaThmDeltaK[Die obere Schrankenmenge liegt im Träger]{%
  T\subseteq A\vdash\UB_{A,\leq}(T)\subseteq A
}{UpperBoundSetSubsetCarrier}{%
  \DeltaRow{Mengen}{A\dsep T\dsep u}
  \DeltaRow{Relationsoperatoren}{\leq}
  \DeltaPrem{Partielle Ordnung}{\PartOrd{A,\leq}}
}
\begin{tabproof}
  \proofstep{1}{T\subseteq A}{\rA}
  \proofstep{2}{u\in\UB_{A,\leq}(T)}{\rA}
  \proofstep{1,2}{u\in A\land\forall x\in T\,(x\leq u)}{%
    \FormulaRefAuto{UpperBoundSetElim}[thm]{1,2}}
  \proofstep{1,2}{u\in A}{\rAEa{3}}
  \proofstep{1}{\UB_{A,\leq}(T)\subseteq A}{%
    \FormulaRefAuto{A \subseteq B \coloneq
      \forall x\,(x\in A \rightarrow x\in B)}{\rUI{\rRI{2,4}}}}
\end{tabproof}

\FormulaThmDeltaK[Charakterisierung der unteren Schrankenmenge]{%
  T\subseteq A\vdash
  \Bigl(u\in\LB_{A,\leq}(T)\leftrightarrow
    \bigl(u\in A\land\forall x\in T\,(u\leq x)\bigr)\Bigr)
}{LowerBoundSetCharacterization}{%
  \DeltaRow{Mengen}{A\dsep T\dsep u\dsep x}
  \DeltaRow{Relationsoperatoren}{\leq}
  \DeltaPrem{Partielle Ordnung}{\PartOrd{A,\leq}}
}
\begin{tabproofwide}
  \proofstepwidestar[1]{T\subseteq A}{\rA}
  \proofstepwidestar[1]{%
    u\in\LB_{A,\leq}(T)\leftrightarrow
    \bigl(u\in A\land\forall x\in T\,(u\leq x)\bigr)}{%
    \FormulaRefAuto{LowerBoundSet}[def]{1}}
\end{tabproofwide}

\FormulaThmDeltaK[Auswertung der Zugehörigkeit zur unteren Schrankenmenge]{%
  T\subseteq A\dsep u\in\LB_{A,\leq}(T)\vdash
  u\in A\land\forall x\in T\,(u\leq x)
}{LowerBoundSetElim}{%
  \DeltaRow{Mengen}{A\dsep T\dsep u\dsep x}
  \DeltaRow{Relationsoperatoren}{\leq}
  \DeltaPrem{Partielle Ordnung}{\PartOrd{A,\leq}}
}
\begin{tabproofwide}
  \proofstepwidestar[1]{T\subseteq A}{\rA}
  \proofstepwidestar[2]{u\in\LB_{A,\leq}(T)}{\rA}
  \proofstepwidestar[1]{%
    u\in\LB_{A,\leq}(T)\leftrightarrow
      \bigl(u\in A\land\forall x\in T\,(u\leq x)\bigr)}{%
    \FormulaRefAuto{LowerBoundSetCharacterization}[thm]{1}}
  \proofstepwidestar[1,2]{u\in A\land\forall x\in T\,(u\leq x)}{%
    \FormulaRefAuto{P\leftrightarrow Q, P\vdash Q}{3,2}}
\end{tabproofwide}

\FormulaThmDeltaK[Einführung in die untere Schrankenmenge]{%
  T\subseteq A\dsep u\in A\dsep
  \forall x\in T\,(u\leq x)\vdash u\in\LB_{A,\leq}(T)
}{LowerBoundSetIntro}{%
  \DeltaRow{Mengen}{A\dsep T\dsep u\dsep x}
  \DeltaRow{Relationsoperatoren}{\leq}
  \DeltaPrem{Partielle Ordnung}{\PartOrd{A,\leq}}
}
\begin{tabproofwide}
  \proofstepwidestar[1]{T\subseteq A}{\rA}
  \proofstepwidestar[2]{u\in A}{\rA}
  \proofstepwidestar[3]{\forall x\in T\,(u\leq x)}{\rA}
  \proofstepwidestar[2,3]{u\in A\land\forall x\in T\,(u\leq x)}{%
    \rAI{2,3}}
  \proofstepwidestar[1]{%
    u\in\LB_{A,\leq}(T)\leftrightarrow
      \bigl(u\in A\land\forall x\in T\,(u\leq x)\bigr)}{%
    \FormulaRefAuto{LowerBoundSetCharacterization}[thm]{1}}
  \proofstepwidestar[1,2,3]{u\in\LB_{A,\leq}(T)}{%
    \FormulaRefAuto{P\leftrightarrow Q, Q\vdash P}{5,4}}
\end{tabproofwide}

\FormulaThmDeltaK[Die untere Schrankenmenge liegt im Träger]{%
  T\subseteq A\vdash\LB_{A,\leq}(T)\subseteq A
}{LowerBoundSetSubsetCarrier}{%
  \DeltaRow{Mengen}{A\dsep T\dsep u}
  \DeltaRow{Relationsoperatoren}{\leq}
  \DeltaPrem{Partielle Ordnung}{\PartOrd{A,\leq}}
}
\begin{tabproof}
  \proofstep{1}{T\subseteq A}{\rA}
  \proofstep{2}{u\in\LB_{A,\leq}(T)}{\rA}
  \proofstep{1,2}{u\in A\land\forall x\in T\,(u\leq x)}{%
    \FormulaRefAuto{LowerBoundSetElim}[thm]{1,2}}
  \proofstep{1,2}{u\in A}{\rAEa{3}}
  \proofstep{1}{\LB_{A,\leq}(T)\subseteq A}{%
    \FormulaRefAuto{A \subseteq B \coloneq
      \forall x\,(x\in A \rightarrow x\in B)}{\rUI{\rRI{2,4}}}}
\end{tabproof}

\FormulaThmDeltaK[Charakterisierung der oberen Schranken eines Elementpaares]{%
  x,y\in A\vdash
  \Bigl(u\in\UB_{A,\leq}(\{x,y\})\leftrightarrow
    \bigl(u\in A\land x\leq u\land y\leq u\bigr)\Bigr)
}{PairUpperBoundSetCharacterization}{%
  \DeltaRow{Mengen}{A\dsep u\dsep x\dsep y\dsep z}
  \DeltaRow{Relationsoperatoren}{\leq}
  \DeltaPrem{Partielle Ordnung}{\PartOrd{A,\leq}}
}
\begin{tabproofsplitwide}
  \proofpartwide{\(\vdash\)}
    \proofstepwidestar[1]{x\in A}{\rA}
    \proofstepwidestar[2]{y\in A}{\rA}
    \proofstepwidestar[3]{u\in\UB_{A,\leq}(\{x,y\})}{\rA}
    \proofstepwidestar[1,2]{\{x,y\}\subseteq A}{%
      \FormulaRefAuto{a \in A,\, b \in A \vdash \{a,b\} \subseteq A}{1,2}}
    \proofstepwidestar[1,2,3]{%
      u\in A\land\forall z\in\{x,y\}\,(z\leq u)}{%
      \FormulaRefAuto{UpperBoundSetElim}[thm]{4,3}}
    \proofstepwidestar[1,2,3]{u\in A}{\rAEa{5}}
    \proofstepwidestar[]{x\in\{x,y\}}{\FormulaRefAuto{a \in \{a,b\}}}
    \proofstepwidestar[]{y\in\{x,y\}}{\FormulaRefAuto{b \in \{a,b\}}}
    \proofstepwidestar[1,2,3]{x\leq u}{\rRE{7,\rUE{\rAEb{5}}}}
    \proofstepwidestar[1,2,3]{y\leq u}{\rRE{8,\rUE{\rAEb{5}}}}
    \proofstepwidestar[1,2,3]{%
      u\in A\land x\leq u\land y\leq u}{\rAI{6,\rAI{9,10}}}
  \closeproofpartwide

  \proofpartwide{\(\dashv\)}
    \proofstepwidestar[1]{x\in A}{\rA}
    \proofstepwidestar[2]{y\in A}{\rA}
    \proofstepwidestar[3]{u\in A\land x\leq u\land y\leq u}{\rA}
    \proofstepwidestar[1,2]{\{x,y\}\subseteq A}{%
      \FormulaRefAuto{a \in A,\, b \in A \vdash \{a,b\} \subseteq A}{1,2}}
    \proofstepwidestar[5]{z\in\{x,y\}}{\rA}
    \proofstepwidestar[5]{z=x\lor z=y}{%
      \FormulaRefAuto{x \in \{A,B\}\;\eqvdash\;(x = A \lor x = B)}{5}}
    \proofstepwidestar[7]{z=x}{\rA}
    \proofstepwidestar[3,7]{z\leq u}{\rIE{7,\rAEa{\rAEb{3}}}}
    \proofstepwidestar[9]{z=y}{\rA}
    \proofstepwidestar[3,9]{z\leq u}{\rIE{9,\rAEb{\rAEb{3}}}}
    \proofstepwidestar[3,5]{z\leq u}{\rOE{6,7,8,9,10}}
    \proofstepwidestar[3]{\forall z\in\{x,y\}\,(z\leq u)}{%
      \rUI{\rRI{5,11}}}
    \proofstepwidestar[1,2,3]{u\in\UB_{A,\leq}(\{x,y\})}{%
      \FormulaRefAuto{UpperBoundSetIntro}[thm]{4,\rAEa{3},12}}
  \closeproofpartwide
\end{tabproofsplitwide}

\FormulaThmDeltaK[Charakterisierung der unteren Schranken eines Elementpaares]{%
  x,y\in A\vdash
  \Bigl(d\in\LB_{A,\leq}(\{x,y\})\leftrightarrow
    \bigl(d\in A\land d\leq x\land d\leq y\bigr)\Bigr)
}{PairLowerBoundSetCharacterization}{%
  \DeltaRow{Mengen}{A\dsep d\dsep x\dsep y\dsep z}
  \DeltaRow{Relationsoperatoren}{\leq}
  \DeltaPrem{Partielle Ordnung}{\PartOrd{A,\leq}}
}
\begin{tabproofsplitwide}
  \proofpartwide{\(\vdash\)}
    \proofstepwidestar[1]{x\in A}{\rA}
    \proofstepwidestar[2]{y\in A}{\rA}
    \proofstepwidestar[3]{d\in\LB_{A,\leq}(\{x,y\})}{\rA}
    \proofstepwidestar[1,2]{\{x,y\}\subseteq A}{%
      \FormulaRefAuto{a \in A,\, b \in A \vdash \{a,b\} \subseteq A}{1,2}}
    \proofstepwidestar[1,2,3]{%
      d\in A\land\forall z\in\{x,y\}\,(d\leq z)}{%
      \FormulaRefAuto{LowerBoundSetElim}[thm]{4,3}}
    \proofstepwidestar[1,2,3]{d\in A}{\rAEa{5}}
    \proofstepwidestar[]{x\in\{x,y\}}{\FormulaRefAuto{a \in \{a,b\}}}
    \proofstepwidestar[]{y\in\{x,y\}}{\FormulaRefAuto{b \in \{a,b\}}}
    \proofstepwidestar[1,2,3]{d\leq x}{\rRE{7,\rUE{\rAEb{5}}}}
    \proofstepwidestar[1,2,3]{d\leq y}{\rRE{8,\rUE{\rAEb{5}}}}
    \proofstepwidestar[1,2,3]{%
      d\in A\land d\leq x\land d\leq y}{\rAI{6,\rAI{9,10}}}
  \closeproofpartwide

  \proofpartwide{\(\dashv\)}
    \proofstepwidestar[1]{x\in A}{\rA}
    \proofstepwidestar[2]{y\in A}{\rA}
    \proofstepwidestar[3]{d\in A\land d\leq x\land d\leq y}{\rA}
    \proofstepwidestar[1,2]{\{x,y\}\subseteq A}{%
      \FormulaRefAuto{a \in A,\, b \in A \vdash \{a,b\} \subseteq A}{1,2}}
    \proofstepwidestar[5]{z\in\{x,y\}}{\rA}
    \proofstepwidestar[5]{z=x\lor z=y}{%
      \FormulaRefAuto{x \in \{A,B\}\;\eqvdash\;(x = A \lor x = B)}{5}}
    \proofstepwidestar[7]{z=x}{\rA}
    \proofstepwidestar[3,7]{d\leq z}{\rIE{7,\rAEa{\rAEb{3}}}}
    \proofstepwidestar[9]{z=y}{\rA}
    \proofstepwidestar[3,9]{d\leq z}{\rIE{9,\rAEb{\rAEb{3}}}}
    \proofstepwidestar[3,5]{d\leq z}{\rOE{6,7,8,9,10}}
    \proofstepwidestar[3]{\forall z\in\{x,y\}\,(d\leq z)}{%
      \rUI{\rRI{5,11}}}
    \proofstepwidestar[1,2,3]{d\in\LB_{A,\leq}(\{x,y\})}{%
      \FormulaRefAuto{LowerBoundSetIntro}[thm]{4,\rAEa{3},12}}
  \closeproofpartwide
\end{tabproofsplitwide}

\FormulaThmDeltaK[Das Supremum ist eine obere Schranke]{%
  \begin{aligned}[t]
    &T\subseteq A\dsep
      \exists m\in\UB_{A,\leq}(T)\,
        \forall a\in\UB_{A,\leq}(T)\,(m\leq a)\vdash{}\\
    &\sup_{A,\leq}(T)\in\UB_{A,\leq}(T)
  \end{aligned}
}{SupremumBelongsUpperBounds}{%
  \DeltaRow{Mengen}{A\dsep T\dsep m\dsep a}
  \DeltaRow{Relationsoperatoren}{\leq}
  \DeltaPrem{Partielle Ordnung}{\PartOrd{A,\leq}}
}
\begin{tabproofwide}
  \proofstepwidestar[1]{T\subseteq A}{\rA}
  \proofstepwidestar[2]{%
    \exists m\in\UB_{A,\leq}(T)\,
      \forall a\in\UB_{A,\leq}(T)\,(m\leq a)}{\rA}
  \proofstepwidestar[2]{%
    \min_{A,\leq}\bigl(\UB_{A,\leq}(T)\bigr)\in\UB_{A,\leq}(T)}{%
    \FormulaRefAuto{Minimum}[def]{2}}
  \proofstepwidestar[1,2]{%
    \sup_{A,\leq}(T)=\min_{A,\leq}\bigl(\UB_{A,\leq}(T)\bigr)}{%
    \FormulaRefAuto{Supremum}[def]{1,2}}
  \proofstepwidestar[1,2]{\sup_{A,\leq}(T)\in\UB_{A,\leq}(T)}{\rIE{4,3}}
\end{tabproofwide}

\FormulaThmDeltaK[Das Supremum liegt im Träger]{%
  \begin{aligned}[t]
    &T\subseteq A\dsep
      \exists m\in\UB_{A,\leq}(T)\,
        \forall a\in\UB_{A,\leq}(T)\,(m\leq a)\vdash{}\\
    &\sup_{A,\leq}(T)\in A
  \end{aligned}
}{SupremumInCarrier}{%
  \DeltaRow{Mengen}{A\dsep T\dsep m\dsep a}
  \DeltaRow{Relationsoperatoren}{\leq}
  \DeltaPrem{Partielle Ordnung}{\PartOrd{A,\leq}}
}
\begin{tabproof}
  \proofstep{1}{T\subseteq A}{\rA}
  \proofstep{2}{%
    \exists m\in\UB_{A,\leq}(T)\,
      \forall a\in\UB_{A,\leq}(T)\,(m\leq a)}{\rA}
  \proofstep{1,2}{\sup_{A,\leq}(T)\in\UB_{A,\leq}(T)}{%
    \FormulaRefAuto{SupremumBelongsUpperBounds}[thm]{1,2}}
  \proofstep{1}{\UB_{A,\leq}(T)\subseteq A}{%
    \FormulaRefAuto{UpperBoundSetSubsetCarrier}[thm]{1}}
  \proofstep{1,2}{\sup_{A,\leq}(T)\in A}{%
    \FormulaRefAuto{A\subseteq B,\,x\in A\vdash x\in B}{4,3}}
\end{tabproof}

\FormulaThmDeltaK[Jedes Element liegt unter dem Supremum]{%
  \begin{aligned}[t]
    &T\subseteq A\dsep
      \exists m\in\UB_{A,\leq}(T)\,
        \forall a\in\UB_{A,\leq}(T)\,(m\leq a)
      \dsep x\in T\vdash{}\\
    &x\leq\sup_{A,\leq}(T)
  \end{aligned}
}{SupremumDominatesMember}{%
  \DeltaRow{Mengen}{A\dsep T\dsep m\dsep a\dsep x}
  \DeltaRow{Relationsoperatoren}{\leq}
  \DeltaPrem{Partielle Ordnung}{\PartOrd{A,\leq}}
}
\begin{tabproof}
  \proofstep{1}{T\subseteq A}{\rA}
  \proofstep{2}{%
    \exists m\in\UB_{A,\leq}(T)\,
      \forall a\in\UB_{A,\leq}(T)\,(m\leq a)}{\rA}
  \proofstep{3}{x\in T}{\rA}
  \proofstep{1,2}{\sup_{A,\leq}(T)\in\UB_{A,\leq}(T)}{%
    \FormulaRefAuto{SupremumBelongsUpperBounds}[thm]{1,2}}
  \proofstep{1,2}{%
    \sup_{A,\leq}(T)\in A\land
      \forall y\in T\,(y\leq\sup_{A,\leq}(T))}{%
    \FormulaRefAuto{UpperBoundSetElim}[thm]{1,4}}
  \proofstep{1,2,3}{x\leq\sup_{A,\leq}(T)}{\rRE{3,\rUE{\rAEb{5}}}}
\end{tabproof}

\FormulaThmDeltaK[Das Supremum ist die kleinste obere Schranke]{%
  \begin{aligned}[t]
    &T\subseteq A\dsep
      \exists m\in\UB_{A,\leq}(T)\,
        \forall b\in\UB_{A,\leq}(T)\,(m\leq b)
      \dsep a\in\UB_{A,\leq}(T)\vdash{}\\
    &\sup_{A,\leq}(T)\leq a
  \end{aligned}
}{SupremumLeastUpperBound}{%
  \DeltaRow{Mengen}{A\dsep T\dsep m\dsep a\dsep b}
  \DeltaRow{Relationsoperatoren}{\leq}
  \DeltaPrem{Partielle Ordnung}{\PartOrd{A,\leq}}
}
\begin{tabproofwide}
  \proofstepwidestar[1]{T\subseteq A}{\rA}
  \proofstepwidestar[2]{%
    \exists m\in\UB_{A,\leq}(T)\,
      \forall b\in\UB_{A,\leq}(T)\,(m\leq b)}{\rA}
  \proofstepwidestar[3]{a\in\UB_{A,\leq}(T)}{\rA}
  \proofstepwidestar[2,3]{%
    \min_{A,\leq}(\UB_{A,\leq}(T))\leq a}{%
    \FormulaRefAuto{Minimum}[def]{2,3}}
  \proofstepwidestar[1,2]{%
    \sup_{A,\leq}(T)=\min_{A,\leq}\bigl(\UB_{A,\leq}(T)\bigr)}{%
    \FormulaRefAuto{Supremum}[def]{1,2}}
  \proofstepwidestar[1,2,3]{\sup_{A,\leq}(T)\leq a}{\rIE{5,4}}
\end{tabproofwide}

\FormulaThmDeltaK[Charakterisierung eines Supremumskandidaten]{%
  \begin{aligned}[t]
    &T\subseteq A\dsep u\in\UB_{A,\leq}(T)\dsep
      \forall a\in\UB_{A,\leq}(T)\,(u\leq a)\vdash{}\\
    &u=\sup_{A,\leq}(T)
  \end{aligned}
}{SupremumCandidateEqualsSupremum}{%
  \DeltaRow{Mengen}{A\dsep T\dsep u\dsep a}
  \DeltaRow{Relationsoperatoren}{\leq}
  \DeltaPrem{Partielle Ordnung}{\PartOrd{A,\leq}}
}
\begin{tabproofwide}
  \proofstepwidestar[1]{T\subseteq A}{\rA}
  \proofstepwidestar[2]{u\in\UB_{A,\leq}(T)}{\rA}
  \proofstepwidestar[3]{\forall a\in\UB_{A,\leq}(T)\,(u\leq a)}{\rA}
  \proofstepwidestar[2,3]{%
    \exists m\in\UB_{A,\leq}(T)\,
      \forall a\in\UB_{A,\leq}(T)\,(m\leq a)}{\rEI{\rAI{2,3}}}
  \proofstepwidestar[1,2,3]{\sup_{A,\leq}(T)\in\UB_{A,\leq}(T)}{%
    \FormulaRefAuto{SupremumBelongsUpperBounds}[thm]{1,4}}
  \proofstepwidestar[1,2,3]{u\leq\sup_{A,\leq}(T)}{\rRE{5,\rUE{3}}}
  \proofstepwidestar[1,2,3]{\sup_{A,\leq}(T)\leq u}{%
    \FormulaRefAuto{SupremumLeastUpperBound}[thm]{1,4,2}}
  \proofstepwidestar[1,2,3]{u=\sup_{A,\leq}(T)}{%
    \FormulaRefAuto{x\leq y \dsep y\leq x \vdash x=y}[ax]{6,7}}
\end{tabproofwide}

\FormulaThmDeltaK[Charakterisierung der Gleichheit mit dem Supremum]{%
  \begin{aligned}[t]
    &T\subseteq A\dsep
      \exists m\in\UB_{A,\leq}(T)\,
        \forall a\in\UB_{A,\leq}(T)\,(m\leq a)
      \dsep u\in A\vdash{}\\
    &u=\sup_{A,\leq}(T)\leftrightarrow
      \Bigl(u\in\UB_{A,\leq}(T)\land
        \forall a\in\UB_{A,\leq}(T)\,(u\leq a)\Bigr)
  \end{aligned}
}{SupremumEqualityCharacterization}{%
  \DeltaRow{Mengen}{A\dsep T\dsep m\dsep u\dsep a}
  \DeltaRow{Relationsoperatoren}{\leq}
  \DeltaPrem{Partielle Ordnung}{\PartOrd{A,\leq}}
}
\begin{tabproofsplitwide}
  \proofpartwideindR[Hinrichtung]{%
    \begin{aligned}[t]
      &T\subseteq A\dsep
        \exists m\in\UB_{A,\leq}(T)\,
          \forall a\in\UB_{A,\leq}(T)\,(m\leq a)
        \dsep u\in A\vdash{}\\
      &u=\sup_{A,\leq}(T)\to
        \Bigl(u\in\UB_{A,\leq}(T)\land
          \forall a\in\UB_{A,\leq}(T)\,(u\leq a)\Bigr)
    \end{aligned}
  }{%
    \begin{aligned}[t]
      &T\subseteq A\dsep
        \exists m\in\UB_{A,\leq}(T)\,
          \forall a\in\UB_{A,\leq}(T)\,(m\leq a)
        \dsep u\in A\vdash{}\\
      &u=\sup_{A,\leq}(T)\to
        \Bigl(u\in\UB_{A,\leq}(T)\land
          \forall a\in\UB_{A,\leq}(T)\,(u\leq a)\Bigr)
    \end{aligned}
  }[SupremumEqualityForward]
    \proofstepwidestar[1]{T\subseteq A}{\rA}
    \proofstepwidestar[2]{%
      \exists m\in\UB_{A,\leq}(T)\,
        \forall a\in\UB_{A,\leq}(T)\,(m\leq a)}{\rA}
    \proofstepwidestar[3]{u\in A}{\rA}
    \proofstepwidestar[4]{u=\sup_{A,\leq}(T)}{\rA}
    \proofstepwidestar[1,2]{\sup_{A,\leq}(T)\in\UB_{A,\leq}(T)}{%
      \FormulaRefAuto{SupremumBelongsUpperBounds}[thm]{1,2}}
    \proofstepwidestar[1,2,4]{u\in\UB_{A,\leq}(T)}{\rIE{4,5}}
    \proofstepwidestar[7]{a\in\UB_{A,\leq}(T)}{\rA}
    \proofstepwidestar[1,2,7]{\sup_{A,\leq}(T)\leq a}{%
      \FormulaRefAuto{SupremumLeastUpperBound}[thm]{1,2,7}}
    \proofstepwidestar[1,2,4,7]{u\leq a}{\rIE{4,8}}
    \proofstepwidestar[1,2,4]{%
      \forall a\in\UB_{A,\leq}(T)\,(u\leq a)}{\rUI{\rRI{7,9}}}
    \proofstepwide[1,2,4]{u\in\UB_{A,\leq}(T)}{\land}{%
      \forall a\in\UB_{A,\leq}(T)\,(u\leq a)}{\rAI{6,10}}
    \proofstepwide[1,2,3]{u=\sup_{A,\leq}(T)}{\to}{%
      \begin{aligned}[t]
        &u\in\UB_{A,\leq}(T)\\[-2pt]
        &{}\land\forall a\in\UB_{A,\leq}(T)\,(u\leq a)
      \end{aligned}}{\rRI{4,11}}
  \closeproofpartwideindR

  \proofpartwideindR[Rückrichtung]{%
    \begin{aligned}[t]
      &T\subseteq A\dsep
        \exists m\in\UB_{A,\leq}(T)\,
          \forall a\in\UB_{A,\leq}(T)\,(m\leq a)
        \dsep u\in A\vdash{}\\
      &\Bigl(u\in\UB_{A,\leq}(T)\land
          \forall a\in\UB_{A,\leq}(T)\,(u\leq a)\Bigr)
        \to u=\sup_{A,\leq}(T)
    \end{aligned}
  }{%
    \begin{aligned}[t]
      &T\subseteq A\dsep
        \exists m\in\UB_{A,\leq}(T)\,
          \forall a\in\UB_{A,\leq}(T)\,(m\leq a)
        \dsep u\in A\vdash{}\\
      &\Bigl(u\in\UB_{A,\leq}(T)\land
          \forall a\in\UB_{A,\leq}(T)\,(u\leq a)\Bigr)
        \to u=\sup_{A,\leq}(T)
    \end{aligned}
  }[SupremumEqualityBackward]
    \proofstepwidestar[1]{T\subseteq A}{\rA}
    \proofstepwidestar[2]{%
      \exists m\in\UB_{A,\leq}(T)\,
        \forall a\in\UB_{A,\leq}(T)\,(m\leq a)}{\rA}
    \proofstepwidestar[3]{u\in A}{\rA}
    \proofstepwidestar[4]{%
      \begin{aligned}[t]
        &u\in\UB_{A,\leq}(T)\\[-2pt]
        &{}\land\forall a\in\UB_{A,\leq}(T)\,(u\leq a)
      \end{aligned}}{\rA}
    \proofstepwidestar[1,4]{u=\sup_{A,\leq}(T)}{%
      \FormulaRefAuto{SupremumCandidateEqualsSupremum}[thm]{%
        1,\rAEa{4},\rAEb{4}}}
    \proofstepwidestar[1,2,3]{%
      \begin{gathered}
        u\in\UB_{A,\leq}(T)\\[-2pt]
        {}\land\forall a\in\UB_{A,\leq}(T)\,(u\leq a)
      \end{gathered}
      \to u=\sup_{A,\leq}(T)}{\rRI{4,5}}
  \closeproofpartwideindR

  \proofpartwide{Schluss}
    \proofstepwidestar[]{%
      u=\sup_{A,\leq}(T)\to
      \Bigl(u\in\UB_{A,\leq}(T)\land
        \forall a\in\UB_{A,\leq}(T)\,(u\leq a)\Bigr)}{%
      \FormulaRefAuto{SupremumEqualityForward}[thm]}
    \proofstepwidestar[]{%
      \Bigl(u\in\UB_{A,\leq}(T)\land
        \forall a\in\UB_{A,\leq}(T)\,(u\leq a)\Bigr)
      \to u=\sup_{A,\leq}(T)}{%
      \FormulaRefAuto{SupremumEqualityBackward}[thm]}
    \proofstepwidestar[]{%
      u=\sup_{A,\leq}(T)\leftrightarrow
      \Bigl(u\in\UB_{A,\leq}(T)\land
        \forall a\in\UB_{A,\leq}(T)\,(u\leq a)\Bigr)}{\rLRI{1,2}}
  \closeproofpartwide
\end{tabproofsplitwide}

\FormulaThmDeltaK[Das Infimum ist eine untere Schranke]{%
  \begin{aligned}[t]
    &T\subseteq A\dsep
      \exists m\in\LB_{A,\leq}(T)\,
        \forall a\in\LB_{A,\leq}(T)\,(a\leq m)\vdash{}\\
    &\inf_{A,\leq}(T)\in\LB_{A,\leq}(T)
  \end{aligned}
}{InfimumBelongsLowerBounds}{%
  \DeltaRow{Mengen}{A\dsep T\dsep m\dsep a}
  \DeltaRow{Relationsoperatoren}{\leq}
  \DeltaPrem{Partielle Ordnung}{\PartOrd{A,\leq}}
}
\begin{tabproofwide}
  \proofstepwidestar[1]{T\subseteq A}{\rA}
  \proofstepwidestar[2]{%
    \exists m\in\LB_{A,\leq}(T)\,
      \forall a\in\LB_{A,\leq}(T)\,(a\leq m)}{\rA}
  \proofstepwidestar[2]{%
    \max_{A,\leq}\bigl(\LB_{A,\leq}(T)\bigr)\in\LB_{A,\leq}(T)}{%
    \FormulaRefAuto{Maximum}[def]{2}}
  \proofstepwidestar[1,2]{%
    \inf_{A,\leq}(T)=\max_{A,\leq}\bigl(\LB_{A,\leq}(T)\bigr)}{%
    \FormulaRefAuto{Infimum}[def]{1,2}}
  \proofstepwidestar[1,2]{\inf_{A,\leq}(T)\in\LB_{A,\leq}(T)}{\rIE{4,3}}
\end{tabproofwide}

\FormulaThmDeltaK[Das Infimum liegt im Träger]{%
  \begin{aligned}[t]
    &T\subseteq A\dsep
      \exists m\in\LB_{A,\leq}(T)\,
        \forall a\in\LB_{A,\leq}(T)\,(a\leq m)\vdash{}\\
    &\inf_{A,\leq}(T)\in A
  \end{aligned}
}{InfimumInCarrier}{%
  \DeltaRow{Mengen}{A\dsep T\dsep m\dsep a}
  \DeltaRow{Relationsoperatoren}{\leq}
  \DeltaPrem{Partielle Ordnung}{\PartOrd{A,\leq}}
}
\begin{tabproof}
  \proofstep{1}{T\subseteq A}{\rA}
  \proofstep{2}{%
    \exists m\in\LB_{A,\leq}(T)\,
      \forall a\in\LB_{A,\leq}(T)\,(a\leq m)}{\rA}
  \proofstep{1,2}{\inf_{A,\leq}(T)\in\LB_{A,\leq}(T)}{%
    \FormulaRefAuto{InfimumBelongsLowerBounds}[thm]{1,2}}
  \proofstep{1}{\LB_{A,\leq}(T)\subseteq A}{%
    \FormulaRefAuto{LowerBoundSetSubsetCarrier}[thm]{1}}
  \proofstep{1,2}{\inf_{A,\leq}(T)\in A}{%
    \FormulaRefAuto{A\subseteq B,\,x\in A\vdash x\in B}{4,3}}
\end{tabproof}

\FormulaThmDeltaK[Das Infimum liegt unter jedem Element]{%
  \begin{aligned}[t]
    &T\subseteq A\dsep
      \exists m\in\LB_{A,\leq}(T)\,
        \forall a\in\LB_{A,\leq}(T)\,(a\leq m)
      \dsep x\in T\vdash{}\\
    &\inf_{A,\leq}(T)\leq x
  \end{aligned}
}{InfimumBelowMember}{%
  \DeltaRow{Mengen}{A\dsep T\dsep m\dsep a\dsep x}
  \DeltaRow{Relationsoperatoren}{\leq}
  \DeltaPrem{Partielle Ordnung}{\PartOrd{A,\leq}}
}
\begin{tabproof}
  \proofstep{1}{T\subseteq A}{\rA}
  \proofstep{2}{%
    \exists m\in\LB_{A,\leq}(T)\,
      \forall a\in\LB_{A,\leq}(T)\,(a\leq m)}{\rA}
  \proofstep{3}{x\in T}{\rA}
  \proofstep{1,2}{\inf_{A,\leq}(T)\in\LB_{A,\leq}(T)}{%
    \FormulaRefAuto{InfimumBelongsLowerBounds}[thm]{1,2}}
  \proofstep{1,2}{%
    \inf_{A,\leq}(T)\in A\land
      \forall y\in T\,(\inf_{A,\leq}(T)\leq y)}{%
    \FormulaRefAuto{LowerBoundSetElim}[thm]{1,4}}
  \proofstep{1,2,3}{\inf_{A,\leq}(T)\leq x}{\rRE{3,\rUE{\rAEb{5}}}}
\end{tabproof}

\FormulaThmDeltaK[Das Infimum ist die größte untere Schranke]{%
  \begin{aligned}[t]
    &T\subseteq A\dsep
      \exists m\in\LB_{A,\leq}(T)\,
        \forall b\in\LB_{A,\leq}(T)\,(b\leq m)
      \dsep a\in\LB_{A,\leq}(T)\vdash{}\\
    &a\leq\inf_{A,\leq}(T)
  \end{aligned}
}{InfimumGreatestLowerBound}{%
  \DeltaRow{Mengen}{A\dsep T\dsep m\dsep a\dsep b}
  \DeltaRow{Relationsoperatoren}{\leq}
  \DeltaPrem{Partielle Ordnung}{\PartOrd{A,\leq}}
}
\begin{tabproofwide}
  \proofstepwidestar[1]{T\subseteq A}{\rA}
  \proofstepwidestar[2]{%
    \exists m\in\LB_{A,\leq}(T)\,
      \forall b\in\LB_{A,\leq}(T)\,(b\leq m)}{\rA}
  \proofstepwidestar[3]{a\in\LB_{A,\leq}(T)}{\rA}
  \proofstepwidestar[2,3]{%
    a\leq\max_{A,\leq}(\LB_{A,\leq}(T))}{%
    \FormulaRefAuto{Maximum}[def]{2,3}}
  \proofstepwidestar[1,2]{%
    \inf_{A,\leq}(T)=\max_{A,\leq}\bigl(\LB_{A,\leq}(T)\bigr)}{%
    \FormulaRefAuto{Infimum}[def]{1,2}}
  \proofstepwidestar[1,2,3]{a\leq\inf_{A,\leq}(T)}{\rIE{5,4}}
\end{tabproofwide}

\FormulaThmDeltaK[Charakterisierung eines Infimumskandidaten]{%
  \begin{aligned}[t]
    &T\subseteq A\dsep u\in\LB_{A,\leq}(T)\dsep
      \forall a\in\LB_{A,\leq}(T)\,(a\leq u)\vdash{}\\
    &u=\inf_{A,\leq}(T)
  \end{aligned}
}{InfimumCandidateEqualsInfimum}{%
  \DeltaRow{Mengen}{A\dsep T\dsep u\dsep a}
  \DeltaRow{Relationsoperatoren}{\leq}
  \DeltaPrem{Partielle Ordnung}{\PartOrd{A,\leq}}
}
\begin{tabproofwide}
  \proofstepwidestar[1]{T\subseteq A}{\rA}
  \proofstepwidestar[2]{u\in\LB_{A,\leq}(T)}{\rA}
  \proofstepwidestar[3]{\forall a\in\LB_{A,\leq}(T)\,(a\leq u)}{\rA}
  \proofstepwidestar[2,3]{%
    \exists m\in\LB_{A,\leq}(T)\,
      \forall a\in\LB_{A,\leq}(T)\,(a\leq m)}{\rEI{\rAI{2,3}}}
  \proofstepwidestar[1,2,3]{\inf_{A,\leq}(T)\in\LB_{A,\leq}(T)}{%
    \FormulaRefAuto{InfimumBelongsLowerBounds}[thm]{1,4}}
  \proofstepwidestar[1,2,3]{\inf_{A,\leq}(T)\leq u}{\rRE{5,\rUE{3}}}
  \proofstepwidestar[1,2,3]{u\leq\inf_{A,\leq}(T)}{%
    \FormulaRefAuto{InfimumGreatestLowerBound}[thm]{1,4,2}}
  \proofstepwidestar[1,2,3]{u=\inf_{A,\leq}(T)}{%
    \FormulaRefAuto{x\leq y \dsep y\leq x \vdash x=y}[ax]{7,6}}
\end{tabproofwide}

\FormulaThmDeltaK[Charakterisierung der Gleichheit mit dem Infimum]{%
  \begin{aligned}[t]
    &T\subseteq A\dsep
      \exists m\in\LB_{A,\leq}(T)\,
        \forall a\in\LB_{A,\leq}(T)\,(a\leq m)
      \dsep u\in A\vdash{}\\
    &u=\inf_{A,\leq}(T)\leftrightarrow
      \Bigl(u\in\LB_{A,\leq}(T)\land
        \forall a\in\LB_{A,\leq}(T)\,(a\leq u)\Bigr)
  \end{aligned}
}{InfimumEqualityCharacterization}{%
  \DeltaRow{Mengen}{A\dsep T\dsep m\dsep u\dsep a}
  \DeltaRow{Relationsoperatoren}{\leq}
  \DeltaPrem{Partielle Ordnung}{\PartOrd{A,\leq}}
}
\begin{tabproofsplitwide}
  \proofpartwideindR[Hinrichtung]{%
    \begin{aligned}[t]
      &T\subseteq A\dsep
        \exists m\in\LB_{A,\leq}(T)\,
          \forall a\in\LB_{A,\leq}(T)\,(a\leq m)
        \dsep u\in A\vdash{}\\
      &u=\inf_{A,\leq}(T)\to
        \Bigl(u\in\LB_{A,\leq}(T)\land
          \forall a\in\LB_{A,\leq}(T)\,(a\leq u)\Bigr)
    \end{aligned}
  }{%
    \begin{aligned}[t]
      &T\subseteq A\dsep
        \exists m\in\LB_{A,\leq}(T)\,
          \forall a\in\LB_{A,\leq}(T)\,(a\leq m)
        \dsep u\in A\vdash{}\\
      &u=\inf_{A,\leq}(T)\to
        \Bigl(u\in\LB_{A,\leq}(T)\land
          \forall a\in\LB_{A,\leq}(T)\,(a\leq u)\Bigr)
    \end{aligned}
  }[InfimumEqualityForward]
    \proofstepwidestar[1]{T\subseteq A}{\rA}
    \proofstepwidestar[2]{%
      \exists m\in\LB_{A,\leq}(T)\,
        \forall a\in\LB_{A,\leq}(T)\,(a\leq m)}{\rA}
    \proofstepwidestar[3]{u\in A}{\rA}
    \proofstepwidestar[4]{u=\inf_{A,\leq}(T)}{\rA}
    \proofstepwidestar[1,2]{\inf_{A,\leq}(T)\in\LB_{A,\leq}(T)}{%
      \FormulaRefAuto{InfimumBelongsLowerBounds}[thm]{1,2}}
    \proofstepwidestar[1,2,4]{u\in\LB_{A,\leq}(T)}{\rIE{4,5}}
    \proofstepwidestar[7]{a\in\LB_{A,\leq}(T)}{\rA}
    \proofstepwidestar[1,2,7]{a\leq\inf_{A,\leq}(T)}{%
      \FormulaRefAuto{InfimumGreatestLowerBound}[thm]{1,2,7}}
    \proofstepwidestar[1,2,4,7]{a\leq u}{\rIE{4,8}}
    \proofstepwidestar[1,2,4]{%
      \forall a\in\LB_{A,\leq}(T)\,(a\leq u)}{\rUI{\rRI{7,9}}}
    \proofstepwide[1,2,4]{u\in\LB_{A,\leq}(T)}{\land}{%
      \forall a\in\LB_{A,\leq}(T)\,(a\leq u)}{\rAI{6,10}}
    \proofstepwide[1,2,3]{u=\inf_{A,\leq}(T)}{\to}{%
      \begin{aligned}[t]
        &u\in\LB_{A,\leq}(T)\\[-2pt]
        &{}\land\forall a\in\LB_{A,\leq}(T)\,(a\leq u)
      \end{aligned}}{\rRI{4,11}}
  \closeproofpartwideindR

  \proofpartwideindR[Rückrichtung]{%
    \begin{aligned}[t]
      &T\subseteq A\dsep
        \exists m\in\LB_{A,\leq}(T)\,
          \forall a\in\LB_{A,\leq}(T)\,(a\leq m)
        \dsep u\in A\vdash{}\\
      &\Bigl(u\in\LB_{A,\leq}(T)\land
          \forall a\in\LB_{A,\leq}(T)\,(a\leq u)\Bigr)
        \to u=\inf_{A,\leq}(T)
    \end{aligned}
  }{%
    \begin{aligned}[t]
      &T\subseteq A\dsep
        \exists m\in\LB_{A,\leq}(T)\,
          \forall a\in\LB_{A,\leq}(T)\,(a\leq m)
        \dsep u\in A\vdash{}\\
      &\Bigl(u\in\LB_{A,\leq}(T)\land
          \forall a\in\LB_{A,\leq}(T)\,(a\leq u)\Bigr)
        \to u=\inf_{A,\leq}(T)
    \end{aligned}
  }[InfimumEqualityBackward]
    \proofstepwidestar[1]{T\subseteq A}{\rA}
    \proofstepwidestar[2]{%
      \exists m\in\LB_{A,\leq}(T)\,
        \forall a\in\LB_{A,\leq}(T)\,(a\leq m)}{\rA}
    \proofstepwidestar[3]{u\in A}{\rA}
    \proofstepwidestar[4]{%
      \begin{aligned}[t]
        &u\in\LB_{A,\leq}(T)\\[-2pt]
        &{}\land\forall a\in\LB_{A,\leq}(T)\,(a\leq u)
      \end{aligned}}{\rA}
    \proofstepwidestar[1,4]{u=\inf_{A,\leq}(T)}{%
      \FormulaRefAuto{InfimumCandidateEqualsInfimum}[thm]{%
        1,\rAEa{4},\rAEb{4}}}
    \proofstepwidestar[1,2,3]{%
      \begin{gathered}
        u\in\LB_{A,\leq}(T)\\[-2pt]
        {}\land\forall a\in\LB_{A,\leq}(T)\,(a\leq u)
      \end{gathered}
      \to u=\inf_{A,\leq}(T)}{\rRI{4,5}}
  \closeproofpartwideindR

  \proofpartwide{Schluss}
    \proofstepwidestar[]{%
      u=\inf_{A,\leq}(T)\to
      \Bigl(u\in\LB_{A,\leq}(T)\land
        \forall a\in\LB_{A,\leq}(T)\,(a\leq u)\Bigr)}{%
      \FormulaRefAuto{InfimumEqualityForward}[thm]}
    \proofstepwidestar[]{%
      \Bigl(u\in\LB_{A,\leq}(T)\land
        \forall a\in\LB_{A,\leq}(T)\,(a\leq u)\Bigr)
      \to u=\inf_{A,\leq}(T)}{%
      \FormulaRefAuto{InfimumEqualityBackward}[thm]}
    \proofstepwidestar[]{%
      u=\inf_{A,\leq}(T)\leftrightarrow
      \Bigl(u\in\LB_{A,\leq}(T)\land
        \forall a\in\LB_{A,\leq}(T)\,(a\leq u)\Bigr)}{\rLRI{1,2}}
  \closeproofpartwide
\end{tabproofsplitwide}

\subsection{Dualität von Schranken und Extremalbegriffen}

Die duale Ordnung vertauscht alle oberen und unteren Begriffe. Damit diese
Aussage nicht in späteren Bänden für einzelne Strukturen wiederholt werden
muss, wird sie hier auf der Ebene beliebiger partieller Ordnungen
festgehalten.

\FormulaThmDeltaK[Untere Schranken der dualen Ordnung]{%
  T\subseteq A\vdash
  \LB_{A,\geq}(T)=\UB_{A,\leq}(T)
}{DualLowerBoundsAreUpperBounds}{%
  \DeltaRow{Mengen}{A\dsep T}
  \DeltaRow{Relationsoperatoren}{\leq\dsep\geq}
  \DeltaPrem{Partielle Ordnung}{\PartOrd{A,\leq}}
  \DeltaPrem{Duale Ordnung}{\PartOrd{A,\geq}}
}

\FormulaThmDeltaK[Obere Schranken der dualen Ordnung]{%
  T\subseteq A\vdash
  \UB_{A,\geq}(T)=\LB_{A,\leq}(T)
}{DualUpperBoundsAreLowerBounds}{%
  \DeltaRow{Mengen}{A\dsep T}
  \DeltaRow{Relationsoperatoren}{\leq\dsep\geq}
  \DeltaPrem{Partielle Ordnung}{\PartOrd{A,\leq}}
  \DeltaPrem{Duale Ordnung}{\PartOrd{A,\geq}}
}

\FormulaThmDeltaK[Minimum und duales Maximum]{%
  T\subseteq A\dsep
  \exists m\in T\,\forall x\in T\,(m\leq x)
  \vdash
  \min_{A,\leq}(T)=\max_{A,\geq}(T)
}{DualMinimumEqualsMaximum}{%
  \DeltaRow{Mengen}{A\dsep T\dsep m\dsep x}
  \DeltaRow{Relationsoperatoren}{\leq\dsep\geq}
  \DeltaPrem{Partielle Ordnung}{\PartOrd{A,\leq}}
}

\FormulaThmDeltaK[Maximum und duales Minimum]{%
  T\subseteq A\dsep
  \exists m\in T\,\forall x\in T\,(x\leq m)
  \vdash
  \max_{A,\leq}(T)=\min_{A,\geq}(T)
}{DualMaximumEqualsMinimum}{%
  \DeltaRow{Mengen}{A\dsep T\dsep m\dsep x}
  \DeltaRow{Relationsoperatoren}{\leq\dsep\geq}
  \DeltaPrem{Partielle Ordnung}{\PartOrd{A,\leq}}
}

\FormulaThmDeltaK[Supremum und duales Infimum]{%
  \begin{aligned}[t]
    &T\subseteq A\dsep
      \exists s\in\UB_{A,\leq}(T)\,
        \forall u\in\UB_{A,\leq}(T)\,(s\leq u)
      \vdash{}\\
    &\sup_{A,\leq}(T)=\inf_{A,\geq}(T)
  \end{aligned}
}{DualSupremumEqualsInfimum}{%
  \DeltaRow{Mengen}{A\dsep T\dsep s\dsep u}
  \DeltaRow{Relationsoperatoren}{\leq\dsep\geq}
  \DeltaPrem{Partielle Ordnung}{\PartOrd{A,\leq}}
}

\FormulaThmDeltaK[Infimum und duales Supremum]{%
  \begin{aligned}[t]
    &T\subseteq A\dsep
      \exists i\in\LB_{A,\leq}(T)\,
        \forall d\in\LB_{A,\leq}(T)\,(d\leq i)
      \vdash{}\\
    &\inf_{A,\leq}(T)=\sup_{A,\geq}(T)
  \end{aligned}
}{DualInfimumEqualsSupremum}{%
  \DeltaRow{Mengen}{A\dsep T\dsep i\dsep d}
  \DeltaRow{Relationsoperatoren}{\leq\dsep\geq}
  \DeltaPrem{Partielle Ordnung}{\PartOrd{A,\leq}}
}

\paragraph{Beweis.}
Für jedes \(a\in A\) gilt nach
\FormulaRefAuto{OrderDualRelation}[def]
\[
  \begin{aligned}
    a\in\LB_{A,\geq}(T)
    &\iff \forall x\in T\,(a\geq x)\\
    &\iff \forall x\in T\,(x\leq a)\\
    &\iff a\in\UB_{A,\leq}(T),\\[2pt]
    a\in\UB_{A,\geq}(T)
    &\iff \forall x\in T\,(x\geq a)\\
    &\iff \forall x\in T\,(a\leq x)\\
    &\iff a\in\LB_{A,\leq}(T).
  \end{aligned}
\]
Die Extensionalität liefert daher die beiden behaupteten Gleichheiten der
Schrankenmengen.

Ein Element \(m\in T\) ist genau dann ein Minimum von \(T\) bezüglich
\(\leq\), wenn \(m\leq x\) für alle \(x\in T\) gilt. Wegen
\(m\leq x\iff x\geq m\) ist dies genau die Bedingung dafür, dass \(m\)
bezüglich \(\geq\) ein Maximum von \(T\) ist. Ebenso ist
\(x\leq m\iff m\geq x\), sodass die Kandidaten für ein Maximum bezüglich
\(\leq\) genau die Kandidaten für ein Minimum bezüglich \(\geq\) sind.
Die Definitionen von Minimum und Maximum ergeben unter den angegebenen
Existenzvoraussetzungen die beiden Gleichheiten.

Mit den bereits bewiesenen Gleichheiten der Schrankenmengen gilt ferner
\[
  \begin{aligned}
    &s\in\UB_{A,\leq}(T)\land
      \forall u\in\UB_{A,\leq}(T)\,(s\leq u)\\
    &\quad\iff
      s\in\LB_{A,\geq}(T)\land
      \forall u\in\LB_{A,\geq}(T)\,(u\geq s),
  \end{aligned}
\]
also ist ein Supremumskandidat bezüglich \(\leq\) genau ein
Infimumskandidat bezüglich \(\geq\). Dual gilt
\[
  \begin{aligned}
    &i\in\LB_{A,\leq}(T)\land
      \forall d\in\LB_{A,\leq}(T)\,(d\leq i)\\
    &\quad\iff
      i\in\UB_{A,\geq}(T)\land
      \forall d\in\UB_{A,\geq}(T)\,(i\geq d).
  \end{aligned}
\]
Die Definitionen von Infimum und Supremum liefern damit unter den jeweiligen
Existenzvoraussetzungen auch die letzten beiden Gleichheiten.
\hfill\(\square\)

\subsection{Beispiele zu Infimum und Supremum}

\subsubsection{Paarmengen}

Für zwei Elemente werden die allgemeinen Auswertungssätze besonders
übersichtlich. In den folgenden Aussagen wird die Existenz des jeweiligen
Paarmengensupremums beziehungsweise Paarmengeninfimums ausdrücklich
vorausgesetzt.

\FormulaDefDeltaK[Existenz eines Paarmengensupremums]{%
  \begin{aligned}[t]
    \operatorname{SupEx}_{A,\leq}(x,y)
    \coloneqq{}&
    x,y\in A\land\\
    &\exists s\in\UB_{A,\leq}(\{x,y\})\,
      \forall u\in\UB_{A,\leq}(\{x,y\})\,(s\leq u)
  \end{aligned}
}{PairSupremumExistence}{%
  \DeltaRow{Mengen}{A\dsep x\dsep y\dsep s\dsep u}
  \DeltaRow{Relationsoperatoren}{\leq}
  \DeltaPrem{Partielle Ordnung}{\PartOrd{A,\leq}}
  \DeltaRow{\textbf{Neues Symbol}}{}
  \DeltaPrem{Existenzprädikat}{\operatorname{SupEx}_{A,\leq}(x,y)}
}

\FormulaDefDeltaK[Existenz eines Paarmengeninfimums]{%
  \begin{aligned}[t]
    \operatorname{InfEx}_{A,\leq}(x,y)
    \coloneqq{}&
    x,y\in A\land\\
    &\exists i\in\LB_{A,\leq}(\{x,y\})\,
      \forall d\in\LB_{A,\leq}(\{x,y\})\,(d\leq i)
  \end{aligned}
}{PairInfimumExistence}{%
  \DeltaRow{Mengen}{A\dsep x\dsep y\dsep i\dsep d}
  \DeltaRow{Relationsoperatoren}{\leq}
  \DeltaPrem{Partielle Ordnung}{\PartOrd{A,\leq}}
  \DeltaRow{\textbf{Neues Symbol}}{}
  \DeltaPrem{Existenzprädikat}{\operatorname{InfEx}_{A,\leq}(x,y)}
}

\FormulaThmDeltaK[Erster Operand liegt unter dem Paarmengensupremum]{%
  \begin{aligned}[t]
    &x,y\in A\dsep
      \exists s\in\UB_{A,\leq}(\{x,y\})\,
        \forall u\in\UB_{A,\leq}(\{x,y\})\,(s\leq u)
      \vdash{}\\
    &x\leq\sup_{A,\leq}(\{x,y\})
  \end{aligned}
}{PairSupremumFirstOperand}{%
  \DeltaRow{Mengen}{A\dsep x\dsep y\dsep s\dsep u}
  \DeltaRow{Relationsoperatoren}{\leq}
  \DeltaPrem{Partielle Ordnung}{\PartOrd{A,\leq}}
}
\begin{tabproofwide}
  \proofstepwidestar[1]{x\in A}{\rA}
  \proofstepwidestar[2]{y\in A}{\rA}
  \proofstepwidestar[3]{%
    \exists s\in\UB_{A,\leq}(\{x,y\})\,
      \forall u\in\UB_{A,\leq}(\{x,y\})\,(s\leq u)}{\rA}
  \proofstepwidestar[1,2]{\{x,y\}\subseteq A}{%
    \FormulaRefAuto{a \in A,\, b \in A \vdash \{a,b\} \subseteq A}{1,2}}
  \proofstepwidestar[]{x\in\{x,y\}}{\FormulaRefAuto{a \in \{a,b\}}}
  \proofstepwidestar[1,2,3]{x\leq\sup_{A,\leq}(\{x,y\})}{%
    \FormulaRefAuto{SupremumDominatesMember}[thm]{4,3,5}}
\end{tabproofwide}

\FormulaThmDeltaK[Zweiter Operand liegt unter dem Paarmengensupremum]{%
  \begin{aligned}[t]
    &x,y\in A\dsep
      \exists s\in\UB_{A,\leq}(\{x,y\})\,
        \forall u\in\UB_{A,\leq}(\{x,y\})\,(s\leq u)
      \vdash{}\\
    &y\leq\sup_{A,\leq}(\{x,y\})
  \end{aligned}
}{PairSupremumSecondOperand}{%
  \DeltaRow{Mengen}{A\dsep x\dsep y\dsep s\dsep u}
  \DeltaRow{Relationsoperatoren}{\leq}
  \DeltaPrem{Partielle Ordnung}{\PartOrd{A,\leq}}
}
\begin{tabproofwide}
  \proofstepwidestar[1]{x\in A}{\rA}
  \proofstepwidestar[2]{y\in A}{\rA}
  \proofstepwidestar[3]{%
    \exists s\in\UB_{A,\leq}(\{x,y\})\,
      \forall u\in\UB_{A,\leq}(\{x,y\})\,(s\leq u)}{\rA}
  \proofstepwidestar[1,2]{\{x,y\}\subseteq A}{%
    \FormulaRefAuto{a \in A,\, b \in A \vdash \{a,b\} \subseteq A}{1,2}}
  \proofstepwidestar[]{y\in\{x,y\}}{\FormulaRefAuto{b \in \{a,b\}}}
  \proofstepwidestar[1,2,3]{y\leq\sup_{A,\leq}(\{x,y\})}{%
    \FormulaRefAuto{SupremumDominatesMember}[thm]{4,3,5}}
\end{tabproofwide}

\FormulaThmDeltaK[Universalcharakterisierung des Paarmengensupremums]{%
  \begin{aligned}[t]
    &x,y,u\in A\dsep{}\\
    &
      \exists s\in\UB_{A,\leq}(\{x,y\})\,
        \forall v\in\UB_{A,\leq}(\{x,y\})\,(s\leq v)\dsep{}\\
    &x\leq u\dsep y\leq u\vdash{}\\
    &\sup_{A,\leq}(\{x,y\})\leq u
  \end{aligned}
}{PairSupremumUniversal}{%
  \DeltaRow{Mengen}{A\dsep x\dsep y\dsep u\dsep s\dsep v}
  \DeltaRow{Relationsoperatoren}{\leq}
  \DeltaPrem{Partielle Ordnung}{\PartOrd{A,\leq}}
}
\begin{tabproofwide}
  \proofstepwidestar[1]{x\in A}{\rA}
  \proofstepwidestar[2]{y\in A}{\rA}
  \proofstepwidestar[3]{u\in A}{\rA}
  \proofstepwidestar[4]{%
    \begin{aligned}[t]
      &\exists s\in\UB_{A,\leq}(\{x,y\})\,\bigl(\\[-2pt]
      &\qquad\forall v\in\UB_{A,\leq}(\{x,y\})\,(s\leq v)\bigr)
    \end{aligned}}{\rA}
  \proofstepwidestar[5]{x\leq u}{\rA}
  \proofstepwidestar[6]{y\leq u}{\rA}
  \proofstepwidestar[1,2]{\{x,y\}\subseteq A}{%
    \FormulaRefAuto{a \in A,\, b \in A \vdash \{a,b\} \subseteq A}{1,2}}
  \proofstepwidestar[5,6]{x\leq u\land y\leq u}{\rAI{5,6}}
  \proofstepwidestar[1,2,3,5,6]{u\in\UB_{A,\leq}(\{x,y\})}{%
    \FormulaRefAuto{P\leftrightarrow Q, Q\vdash P}{%
      \FormulaRefAuto{PairUpperBoundSetCharacterization}[thm]{1,2},
      \rAI{3,8}}}
  \proofstepwidestar[1,2,3,4,5,6]{%
    \sup_{A,\leq}(\{x,y\})\leq u}{%
    \FormulaRefAuto{SupremumLeastUpperBound}[thm]{7,4,9}}
\end{tabproofwide}

\FormulaThmDeltaK[Erster Operand liegt über dem Paarmengeninfimum]{%
  \begin{aligned}[t]
    &x,y\in A\dsep
      \exists i\in\LB_{A,\leq}(\{x,y\})\,
        \forall d\in\LB_{A,\leq}(\{x,y\})\,(d\leq i)
      \vdash{}\\
    &\inf_{A,\leq}(\{x,y\})\leq x
  \end{aligned}
}{PairInfimumFirstOperand}{%
  \DeltaRow{Mengen}{A\dsep x\dsep y\dsep i\dsep d}
  \DeltaRow{Relationsoperatoren}{\leq}
  \DeltaPrem{Partielle Ordnung}{\PartOrd{A,\leq}}
}
\begin{tabproofwide}
  \proofstepwidestar[1]{x\in A}{\rA}
  \proofstepwidestar[2]{y\in A}{\rA}
  \proofstepwidestar[3]{%
    \exists i\in\LB_{A,\leq}(\{x,y\})\,
      \forall d\in\LB_{A,\leq}(\{x,y\})\,(d\leq i)}{\rA}
  \proofstepwidestar[1,2]{\{x,y\}\subseteq A}{%
    \FormulaRefAuto{a \in A,\, b \in A \vdash \{a,b\} \subseteq A}{1,2}}
  \proofstepwidestar[]{x\in\{x,y\}}{\FormulaRefAuto{a \in \{a,b\}}}
  \proofstepwidestar[1,2,3]{\inf_{A,\leq}(\{x,y\})\leq x}{%
    \FormulaRefAuto{InfimumBelowMember}[thm]{4,3,5}}
\end{tabproofwide}

\FormulaThmDeltaK[Zweiter Operand liegt über dem Paarmengeninfimum]{%
  \begin{aligned}[t]
    &x,y\in A\dsep
      \exists i\in\LB_{A,\leq}(\{x,y\})\,
        \forall d\in\LB_{A,\leq}(\{x,y\})\,(d\leq i)
      \vdash{}\\
    &\inf_{A,\leq}(\{x,y\})\leq y
  \end{aligned}
}{PairInfimumSecondOperand}{%
  \DeltaRow{Mengen}{A\dsep x\dsep y\dsep i\dsep d}
  \DeltaRow{Relationsoperatoren}{\leq}
  \DeltaPrem{Partielle Ordnung}{\PartOrd{A,\leq}}
}
\begin{tabproofwide}
  \proofstepwidestar[1]{x\in A}{\rA}
  \proofstepwidestar[2]{y\in A}{\rA}
  \proofstepwidestar[3]{%
    \exists i\in\LB_{A,\leq}(\{x,y\})\,
      \forall d\in\LB_{A,\leq}(\{x,y\})\,(d\leq i)}{\rA}
  \proofstepwidestar[1,2]{\{x,y\}\subseteq A}{%
    \FormulaRefAuto{a \in A,\, b \in A \vdash \{a,b\} \subseteq A}{1,2}}
  \proofstepwidestar[]{y\in\{x,y\}}{\FormulaRefAuto{b \in \{a,b\}}}
  \proofstepwidestar[1,2,3]{\inf_{A,\leq}(\{x,y\})\leq y}{%
    \FormulaRefAuto{InfimumBelowMember}[thm]{4,3,5}}
\end{tabproofwide}

\FormulaThmDeltaK[Universalcharakterisierung des Paarmengeninfimums]{%
  \begin{aligned}[t]
    &x,y,d\in A\dsep{}\\
    &
      \exists i\in\LB_{A,\leq}(\{x,y\})\,
        \forall e\in\LB_{A,\leq}(\{x,y\})\,(e\leq i)\dsep{}\\
    &d\leq x\dsep d\leq y\vdash{}\\
    &d\leq\inf_{A,\leq}(\{x,y\})
  \end{aligned}
}{PairInfimumUniversal}{%
  \DeltaRow{Mengen}{A\dsep x\dsep y\dsep d\dsep i\dsep e}
  \DeltaRow{Relationsoperatoren}{\leq}
  \DeltaPrem{Partielle Ordnung}{\PartOrd{A,\leq}}
}
\begin{tabproofwide}
  \proofstepwidestar[1]{x\in A}{\rA}
  \proofstepwidestar[2]{y\in A}{\rA}
  \proofstepwidestar[3]{d\in A}{\rA}
  \proofstepwidestar[4]{%
    \begin{aligned}[t]
      &\exists i\in\LB_{A,\leq}(\{x,y\})\,\bigl(\\[-2pt]
      &\qquad\forall e\in\LB_{A,\leq}(\{x,y\})\,(e\leq i)\bigr)
    \end{aligned}}{\rA}
  \proofstepwidestar[5]{d\leq x}{\rA}
  \proofstepwidestar[6]{d\leq y}{\rA}
  \proofstepwidestar[1,2]{\{x,y\}\subseteq A}{%
    \FormulaRefAuto{a \in A,\, b \in A \vdash \{a,b\} \subseteq A}{1,2}}
  \proofstepwidestar[5,6]{d\leq x\land d\leq y}{\rAI{5,6}}
  \proofstepwidestar[1,2,3,5,6]{d\in\LB_{A,\leq}(\{x,y\})}{%
    \FormulaRefAuto{P\leftrightarrow Q, Q\vdash P}{%
      \FormulaRefAuto{PairLowerBoundSetCharacterization}[thm]{1,2},
      \rAI{3,8}}}
  \proofstepwidestar[1,2,3,4,5,6]{%
    d\leq\inf_{A,\leq}(\{x,y\})}{%
    \FormulaRefAuto{InfimumGreatestLowerBound}[thm]{7,4,9}}
\end{tabproofwide}

\FormulaThmDeltaK[Idempotenz des Paarmengensupremums]{%
  x\in A\vdash\sup_{A,\leq}(\{x,x\})=x
}{PairSupremumIdempotent}{%
  \DeltaRow{Mengen}{A\dsep x}
  \DeltaRow{Relationsoperatoren}{\leq}
  \DeltaPrem{Partielle Ordnung}{\PartOrd{A,\leq}}
}

\FormulaThmDeltaK[Idempotenz des Paarmengeninfimums]{%
  x\in A\vdash\inf_{A,\leq}(\{x,x\})=x
}{PairInfimumIdempotent}{%
  \DeltaRow{Mengen}{A\dsep x}
  \DeltaRow{Relationsoperatoren}{\leq}
  \DeltaPrem{Partielle Ordnung}{\PartOrd{A,\leq}}
}

\FormulaThmDeltaK[Kommutativität des Paarmengensupremums]{%
  \begin{aligned}[t]
    &x,y\in A\dsep
      \exists s\in\UB_{A,\leq}(\{x,y\})\,
        \forall u\in\UB_{A,\leq}(\{x,y\})\,(s\leq u)
      \vdash{}\\
    &\sup_{A,\leq}(\{x,y\})=\sup_{A,\leq}(\{y,x\})
  \end{aligned}
}{PairSupremumCommutative}{%
  \DeltaRow{Mengen}{A\dsep x\dsep y\dsep s\dsep u}
  \DeltaRow{Relationsoperatoren}{\leq}
  \DeltaPrem{Partielle Ordnung}{\PartOrd{A,\leq}}
}

\FormulaThmDeltaK[Kommutativität des Paarmengeninfimums]{%
  \begin{aligned}[t]
    &x,y\in A\dsep
      \exists i\in\LB_{A,\leq}(\{x,y\})\,
        \forall d\in\LB_{A,\leq}(\{x,y\})\,(d\leq i)
      \vdash{}\\
    &\inf_{A,\leq}(\{x,y\})=\inf_{A,\leq}(\{y,x\})
  \end{aligned}
}{PairInfimumCommutative}{%
  \DeltaRow{Mengen}{A\dsep x\dsep y\dsep i\dsep d}
  \DeltaRow{Relationsoperatoren}{\leq}
  \DeltaPrem{Partielle Ordnung}{\PartOrd{A,\leq}}
}

\FormulaThmDeltaK[Assoziativität des Paarmengensupremums]{%
  \begin{aligned}[t]
    &\operatorname{SupEx}_{A,\leq}(x,y)
      \dsep\operatorname{SupEx}_{A,\leq}(y,z)\\
    &\dsep\operatorname{SupEx}_{A,\leq}
      \bigl(\sup_{A,\leq}(\{x,y\}),z\bigr)
      \dsep\operatorname{SupEx}_{A,\leq}
      \bigl(x,\sup_{A,\leq}(\{y,z\})\bigr)\vdash{}\\
    &\sup_{A,\leq}
       \bigl(\{\sup_{A,\leq}(\{x,y\}),z\}\bigr)
     =
     \sup_{A,\leq}
       \bigl(\{x,\sup_{A,\leq}(\{y,z\})\}\bigr)
  \end{aligned}
}{PairSupremumAssociative}{%
  \DeltaRow{Mengen}{A\dsep x\dsep y\dsep z}
  \DeltaRow{Relationsoperatoren}{\leq}
  \DeltaPrem{Partielle Ordnung}{\PartOrd{A,\leq}}
}

\FormulaThmDeltaK[Assoziativität des Paarmengeninfimums]{%
  \begin{aligned}[t]
    &\operatorname{InfEx}_{A,\leq}(x,y)
      \dsep\operatorname{InfEx}_{A,\leq}(y,z)\\
    &\dsep\operatorname{InfEx}_{A,\leq}
      \bigl(\inf_{A,\leq}(\{x,y\}),z\bigr)
      \dsep\operatorname{InfEx}_{A,\leq}
      \bigl(x,\inf_{A,\leq}(\{y,z\})\bigr)\vdash{}\\
    &\inf_{A,\leq}
       \bigl(\{\inf_{A,\leq}(\{x,y\}),z\}\bigr)
     =
     \inf_{A,\leq}
       \bigl(\{x,\inf_{A,\leq}(\{y,z\})\}\bigr)
  \end{aligned}
}{PairInfimumAssociative}{%
  \DeltaRow{Mengen}{A\dsep x\dsep y\dsep z}
  \DeltaRow{Relationsoperatoren}{\leq}
  \DeltaPrem{Partielle Ordnung}{\PartOrd{A,\leq}}
}

\paragraph{Beweis.}
Es ist \(\{x,x\}=\{x\}\). Wegen der Reflexivität ist \(x\) eine obere
Schranke dieser Menge, und jede obere Schranke \(u\) erfüllt \(x\leq u\).
Somit ist \(x\) der kleinste Kandidat in
\(\UB_{A,\leq}(\{x\})\). Dual ist \(x\) eine untere Schranke, und jede
untere Schranke \(d\) erfüllt \(d\leq x\); daher ist \(x\) der größte
Kandidat in \(\LB_{A,\leq}(\{x\})\). Nach den Definitionen von Supremum
und Infimum folgen beide Idempotenzgleichungen.

Aus \(\{x,y\}=\{y,x\}\) folgen unmittelbar
\[
  \UB_{A,\leq}(\{x,y\})=\UB_{A,\leq}(\{y,x\})
  \quad\text{und}\quad
  \LB_{A,\leq}(\{x,y\})=\LB_{A,\leq}(\{y,x\}).
\]
Die jeweils betrachteten Mengen besitzen folglich denselben eindeutig
bestimmten kleinsten beziehungsweise größten Kandidaten. Die Definitionen
von Supremum und Infimum ergeben die beiden Kommutativitätsgleichungen.

Für den assoziativen Supremumsfall setze
\[
  \begin{aligned}
    p&\coloneqq
      \sup_{A,\leq}\bigl(\{\sup_{A,\leq}(\{x,y\}),z\}\bigr),\\
    q&\coloneqq
      \sup_{A,\leq}\bigl(\{x,\sup_{A,\leq}(\{y,z\})\}\bigr).
  \end{aligned}
\]
Die Existenzvoraussetzungen sichern die Existenz aller hierbei verwendeten
Suprema. Die beiden Operandenregeln, die Transitivität und die
Universalcharakterisierung des Paarmengensupremums liefern für jedes
\(u\in A\)
\[
  \begin{aligned}
    p\leq u
    &\iff x\leq u\land y\leq u\land z\leq u\\
    &\iff q\leq u.
  \end{aligned}
\]
Damit sind \(p\) und \(q\) beide die kleinste obere Schranke von
\(\{x,y,z\}\). Insbesondere gelten \(p\leq q\) und \(q\leq p\), also
folgt \(p=q\) aus der Antisymmetrie.

Für den dualen Infimumsfall setze
\[
  \begin{aligned}
    r&\coloneqq
      \inf_{A,\leq}\bigl(\{\inf_{A,\leq}(\{x,y\}),z\}\bigr),\\
    t&\coloneqq
      \inf_{A,\leq}\bigl(\{x,\inf_{A,\leq}(\{y,z\})\}\bigr).
  \end{aligned}
\]
Die entsprechenden Existenzvoraussetzungen und die dualen Operandenregeln
sowie die Universalcharakterisierung des Paarmengeninfimums ergeben für
jedes \(d\in A\)
\[
  \begin{aligned}
    d\leq r
    &\iff d\leq x\land d\leq y\land d\leq z\\
    &\iff d\leq t.
  \end{aligned}
\]
Also sind \(r\) und \(t\) beide die größte untere Schranke von
\(\{x,y,z\}\). Es folgen \(r\leq t\) und \(t\leq r\), mithin \(r=t\)
durch Antisymmetrie.
\hfill\(\square\)

\subsubsection{Inklusionsordnung auf Mengenfamilien}

\FormulaThmDeltaK[Inklusionsordnung auf einer Mengenfamilie]{%
  M\subseteq\powerset(U)\vdash\PartOrd{M,\subseteq}
}{InclusionOrderOnFamily}{%
  \DeltaRow{Mengen}{U\dsep M\dsep X\dsep Y\dsep Z}
  \DeltaRow{Relationsoperatoren}{\subseteq}
}
\begin{tabproofwide}
  \proofstepwidestar[1]{M\subseteq\powerset(U)}{\rA}
  \proofstepwidestar[]{X\in M\vdash X\subseteq X}{%
    \FormulaRefAuto{A \subseteq A}}
  \proofstepwidestar[]{%
    X\subseteq Y\dsep Y\subseteq Z\vdash X\subseteq Z}{%
    \FormulaRefAuto{A \subseteq B, B \subseteq C \vdash A \subseteq C}}
  \proofstepwidestar[]{%
    X\subseteq Y\dsep Y\subseteq X\vdash X=Y}{%
    \FormulaRefAuto{A \subseteq B, B \subseteq A \vdash A = B}}
  \proofstepwidestar[1]{\PartOrd{M,\subseteq}}{%
    \FormulaRefAuto{Partielle Ordnung}[def]{2,3,4}}
\end{tabproofwide}

Die Reflexivität, Transitivität und Antisymmetrie dieser Ordnung sind genau
die entsprechenden Grundgesetze der Teilmengenrelation.

\FormulaThmDeltaK[Paarmengeninfimum in einer Mengenfamilie]{%
  \begin{aligned}[t]
    &M\subseteq\powerset(U)\dsep
      X,Y\in M\dsep X\cap Y\in M\vdash{}\\
    &\inf_{M,\subseteq}(\{X,Y\})=X\cap Y
  \end{aligned}
}{InclusionFamilyPairInfimum}{%
  \DeltaRow{Mengen}{U\dsep M\dsep X\dsep Y}
  \DeltaRow{Relationsoperatoren}{\subseteq}
  \DeltaPrem{Partielle Ordnung}{\PartOrd{M,\subseteq}}
}

\FormulaThmDeltaK[Paarmengensupremum in einer Mengenfamilie]{%
  \begin{aligned}[t]
    &M\subseteq\powerset(U)\dsep
      X,Y\in M\dsep X\cup Y\in M\vdash{}\\
    &\sup_{M,\subseteq}(\{X,Y\})=X\cup Y
  \end{aligned}
}{InclusionFamilyPairSupremum}{%
  \DeltaRow{Mengen}{U\dsep M\dsep X\dsep Y}
  \DeltaRow{Relationsoperatoren}{\subseteq}
  \DeltaPrem{Partielle Ordnung}{\PartOrd{M,\subseteq}}
}

Tatsächlich ist \(X\cap Y\) in der Inklusionsordnung eine untere Schranke
von \(X\) und \(Y\), und jede weitere untere Schranke ist in
\(X\cap Y\) enthalten. Dual dazu enthält \(X\cup Y\) beide Mengen, und jede
gemeinsame obere Schranke enthält \(X\cup Y\). Die Zusatzannahmen
\(X\cap Y\in M\) beziehungsweise \(X\cup Y\in M\) sind wesentlich: Ein
Infimum oder Supremum muss im gewählten Träger liegen.

\FormulaThmDeltaK[Paarmengeninfimum im Potenzmengenträger]{%
  X,Y\subseteq U\vdash
  \inf_{\powerset(U),\subseteq}(\{X,Y\})=X\cap Y
}{PowerSetPairInfimumIntersection}{%
  \DeltaRow{Mengen}{U\dsep X\dsep Y}
  \DeltaRow{Relationsoperatoren}{\subseteq}
}

\FormulaThmDeltaK[Paarmengensupremum im Potenzmengenträger]{%
  X,Y\subseteq U\vdash
  \sup_{\powerset(U),\subseteq}(\{X,Y\})=X\cup Y
}{PowerSetPairSupremumUnion}{%
  \DeltaRow{Mengen}{U\dsep X\dsep Y}
  \DeltaRow{Relationsoperatoren}{\subseteq}
}

Der Potenzmengenträger ist unter Schnitt und Vereinigung abgeschlossen.
Daher entfallen dort die zusätzlichen Trägerannahmen der beiden
vorangegangenen Sätze.


\subsection{Restriktion einer partiellen Ordnung}

\FormulaDefDeltaK[Restriktion eines Ordnungsoperators]{%
  x,y\in T\vdash
  \bigl(x\leq_T y\leftrightarrow x\leq y\bigr)
}{OrderRestrictionOperator}{%
  \DeltaRow{Mengen}{A\dsep T\dsep x\dsep y}
  \DeltaRow{Relationsoperatoren}{\leq\dsep\leq_T}
  \DeltaPrem{Teilmenge}{T\subseteq A}
  \DeltaRow{\textbf{Neues Symbol}}{}
  \DeltaPrem{Restriktionsoperator}{\leq_T}
}

Der Index \(T\) macht den neuen Träger sichtbar. Die Restriktion ändert
also nicht die Vergleichsaussage, sondern nur den Bereich, auf dem der
Relationsoperator ausgewertet wird.

\FormulaThmDeltaK[Reflexivität der Restriktion]{%
  T\subseteq A\dsep x\in T\vdash x\leq_T x
}{OrderRestrictionReflexive}{%
  \DeltaRow{Mengen}{A\dsep T\dsep x}
  \DeltaRow{Relationsoperatoren}{\leq\dsep\leq_T}
  \DeltaPrem{Partielle Ordnung}{\PartOrd{A,\leq}}
}
\begin{tabproof}
  \proofstep{1}{T\subseteq A}{\rA}
  \proofstep{2}{x\in T}{\rA}
  \proofstep{1,2}{x\in A}{%
    \FormulaRefAuto{A\subseteq B, x\in A \vdash x\in B}{1,2}}
  \proofstep{1,2}{x\leq x}{%
    \FormulaRefAuto{EqRelReflexivity}[ax]{3}}
  \proofstep{2}{x\leq_T x\leftrightarrow x\leq x}{%
    \FormulaRefAuto{OrderRestrictionOperator}[def]{2,2}}
  \proofstep{1,2}{x\leq_T x}{%
    \FormulaRefAuto{P\leftrightarrow Q, Q\vdash P}{5,4}}
\end{tabproof}

\FormulaThmDeltaK[Transitivität der Restriktion]{%
  x,y,z\in T\dsep x\leq_T y\dsep y\leq_T z\vdash x\leq_T z
}{OrderRestrictionTransitive}{%
  \DeltaRow{Mengen}{A\dsep T\dsep x\dsep y\dsep z}
  \DeltaRow{Relationsoperatoren}{\leq\dsep\leq_T}
  \DeltaPrem{Teilmenge}{T\subseteq A}
  \DeltaPrem{Partielle Ordnung}{\PartOrd{A,\leq}}
}
\begin{tabproofwide}
  \proofstepwidestar[1]{x\in T}{\rA}
  \proofstepwidestar[2]{y\in T}{\rA}
  \proofstepwidestar[3]{z\in T}{\rA}
  \proofstepwidestar[4]{x\leq_T y}{\rA}
  \proofstepwidestar[5]{y\leq_T z}{\rA}
  \proofstepwidestar[1,2,4]{x\leq y}{%
    \FormulaRefAuto{P\leftrightarrow Q, P\vdash Q}{%
      \FormulaRefAuto{OrderRestrictionOperator}[def]{1,2},4}}
  \proofstepwidestar[2,3,5]{y\leq z}{%
    \FormulaRefAuto{P\leftrightarrow Q, P\vdash Q}{%
      \FormulaRefAuto{OrderRestrictionOperator}[def]{2,3},5}}
  \proofstepwidestar[1,2,3,4,5]{x\leq z}{%
    \FormulaRefAuto{EqRelTransitivity}[ax]{6,7}}
  \proofstepwidestar[1,3]{x\leq_T z\leftrightarrow x\leq z}{%
    \FormulaRefAuto{OrderRestrictionOperator}[def]{1,3}}
  \proofstepwidestar[1,2,3,4,5]{x\leq_T z}{%
    \FormulaRefAuto{P\leftrightarrow Q, Q\vdash P}{9,8}}
\end{tabproofwide}


\FormulaThmDeltaK[Antisymmetrie der Restriktion]{%
  x,y\in T\dsep x\leq_T y\dsep y\leq_T x\vdash x=y
}{OrderRestrictionAntisymmetric}{%
  \DeltaRow{Mengen}{A\dsep T\dsep x\dsep y}
  \DeltaRow{Relationsoperatoren}{\leq\dsep\leq_T}
  \DeltaPrem{Teilmenge}{T\subseteq A}
  \DeltaPrem{Partielle Ordnung}{\PartOrd{A,\leq}}
}
\begin{tabproofwide}
  \proofstepwidestar[1]{x\in T}{\rA}
  \proofstepwidestar[2]{y\in T}{\rA}
  \proofstepwidestar[3]{x\leq_T y}{\rA}
  \proofstepwidestar[4]{y\leq_T x}{\rA}
  \proofstepwidestar[1,2,3]{x\leq y}{%
    \FormulaRefAuto{P\leftrightarrow Q, P\vdash Q}{%
      \FormulaRefAuto{OrderRestrictionOperator}[def]{1,2},3}}
  \proofstepwidestar[1,2,4]{y\leq x}{%
    \FormulaRefAuto{P\leftrightarrow Q, P\vdash Q}{%
      \FormulaRefAuto{OrderRestrictionOperator}[def]{2,1},4}}
  \proofstepwidestar[1,2,3,4]{x=y}{%
    \FormulaRefAuto{x\leq y \dsep y\leq x \vdash x=y}[ax]{5,6}}
\end{tabproofwide}


\FormulaThmDeltaK[Restriktion einer partiellen Ordnung]{%
  T\subseteq A\dsep\PartOrd{A,\leq}\vdash\PartOrd{T,\leq_T}
}{OrderRestrictionPartialOrder}{%
  \DeltaRow{Mengen}{A\dsep T\dsep x\dsep y\dsep z}
  \DeltaRow{Relationsoperatoren}{\leq\dsep\leq_T}
  \DeltaPrem{Teilmenge}{T\subseteq A}
  \DeltaPrem{Partielle Ordnung}{\PartOrd{A,\leq}}
}
\begin{tabproofwide}
  \proofstepwidestar[1]{T\subseteq A}{\rA}
  \proofstepwidestar[2]{\PartOrd{A,\leq}}{\rA}
  \proofstepwidestar[1,2]{x\in T\vdash x\leq_T x}{%
    \FormulaRefAuto{OrderRestrictionReflexive}[thm]{1}}
  \proofstepwidestar[1,2]{%
    x,y,z\in T\dsep x\leq_T y\dsep y\leq_T z\vdash x\leq_T z}{%
    \FormulaRefAuto{OrderRestrictionTransitive}[thm]}
  \proofstepwidestar[1,2]{%
    x,y\in T\dsep x\leq_T y\dsep y\leq_T x\vdash x=y}{%
    \FormulaRefAuto{OrderRestrictionAntisymmetric}[thm]}
  \proofstepwidestar[1,2]{\PartOrd{T,\leq_T}}{%
    \FormulaRefAuto{Partielle Ordnung}[def]{3,4,5}}
\end{tabproofwide}



% ============================================================
% Wohlordnung
% (totale Ordnung + jedes nichtleere T ⊆ A besitzt ein Minimum)
% ============================================================

\subsection{Wohlordnung}
\subsubsection{Axiom einer Wohlordnung}

\FormulaAxiomDelta[Wohlordnungsprinzip]{%
  T\subseteq A \dsep T\neq\varnothing
  \vdash \exists m\in T\,\forall x\in T\,\bigl(m\leq x\bigr)%
}{
  \DeltaRow{Mengen}{A\dsep T \dsep m \dsep x}
  \DeltaRow{Relationsoperatoren}{\leq}
}

\subsubsection{Definition der Wohlordnung}

\FormulaDefDeltaK[Begriff der Wohlordnung]{\WellOrd{A,\leq}}{Wohlordnung}{
  \DeltaRow{Mengen}{A\dsep T \dsep m \dsep x \dsep y \dsep z}
  \DeltaRow{Relationsoperatoren}{\leq}
  \DeltaRow{\textbf{Axiome}}{}
  \DeltaRow{Totale Ordnung}
           {\TotOrd{A,\leq}}
           [\FormulaRefAuto{Totale Ordnung}]
  \DeltaRow{Wohlordnungsprinzip}{%
    \begin{aligned}
      &T\subseteq A \dsep T\neq\varnothing\\
      &\vdash \exists m\in T\,\forall x\in T\,\bigl(m\leq x\bigr)
    \end{aligned}}
           [\FormulaRefAuto{T\subseteq A \dsep T\neq\varnothing \vdash \exists m\in T\,\forall x\in T\,\bigl(m\leq x\bigr)}]
  \DeltaRow{\textbf{Neue Symbole}}{}
  \DeltaPrem{Wohlordnung}{\WellOrd{A,\leq}}
}

\FormulaDefDeltaK[Begriff der Wohlordenbarkeit]{%
  \WellOrdable{A}\coloneqq
  \exists\,\preccurlyeq\;\WellOrd{A,\preccurlyeq}%
}{WellOrdable}{
  \DeltaRow{Mengen}{A}
  \DeltaRow{Relationsoperatoren}{\preccurlyeq}
  \DeltaRow{Wohlordnung}
           {\WellOrd{A,\preccurlyeq}}
           [\FormulaRefAuto{Wohlordnung}]
  \DeltaRow{\textbf{Neue Symbole}}{}
  \DeltaPrem{Wohlordenbare Menge}{\WellOrdable{A}}
}

\FormulaDefDeltaK[Wohlordnungssatz]{%
  \WOP\coloneqq \forall A\,\WellOrdable{A}%
}{WOP}{
  \DeltaRow{Mengen}{A}
  \DeltaRow{Wohlordenbarkeit}
           {\WellOrdable{A}}
           [\FormulaRefAuto{WellOrdable}]
}





\subsubsection{Wohlordnungssatz}

\FormulaDefDelta[Auswahlrelation auf nichtleeren Teilmengen]{%
  \ChoiceRel{A}\coloneqq
  \{(X,a)\in \PowNE{A}\times A\mid a\in X\}%
}{
  \DeltaDecl{Mengen}{A \dsep X \dsep a}
  \DeltaDecl{Relation}{\ChoiceRel{A}}
}

\FormulaDefDeltaK[Auswahlfunktion auf nichtleeren Teilmengen]{%
  \ChoicePow{C}{A}\coloneqq
  C\colon \PowNE{A}\to A
  \land \forall X\in \PowNE{A}\,\bigl(C(X)\in X\bigr)%
}{ChoicePow}{
  \DeltaDecl{Mengen}{A \dsep X}
  \DeltaDecl{Funktion}{C}
  \DeltaRow{Definitionsbereich}{\PowNE{A}}
  \DeltaRow{Auswahlbedingung}
           {\forall X\in \PowNE{A}\,\bigl(C(X)\in X\bigr)}
}

\FormulaThmDeltaK[Auswahlrelation ist total]{%
  \TotRel{\ChoiceRel{A},\PowNE{A},A}%
}{ChoiceRelTotal}{
  \DeltaDecl{Mengen}{A \dsep X \dsep a}
  \DeltaDecl{Relation}{\ChoiceRel{A}}
}
\begin{tabproof}
  \proofstep{}{%
    \ChoiceRel{A}\subseteq \PowNE{A}\times A%
  }{%
    \FormulaRefAuto{A=\{ x \in B \mid P(x)\}\vdash  A\subseteq B}{%
      \FormulaRefAuto{\ChoiceRel{A}\coloneqq
      \{(X,a)\in \PowNE{A}\times A\mid a\in X\}}}%
  }

  \proofstep{2}{X\in \PowNE{A}}{\rA}
  \proofstep{2}{X\in \powerset(A)}{%
    \FormulaRefAuto{x \in A \setminus B \vdash x \in A}{2}}
  \proofstep{2}{X\notin \{\varnothing\}}{%
    \FormulaRefAuto{x\in A\setminus B\vdash x\notin B}{2}}
  \proofstep{2}{X\neq\varnothing}{%
    \FormulaRefAuto{x \notin \{a\} \eqvdash x \neq a}{4}}
  \proofstep{2}{X\subseteq A}{%
    \FormulaRefAuto{x \in \powerset(A)\eqvdash x \subseteq A}{3}}
  \proofstep{2}{\exists a\,(a\in X)}{%
    \FormulaRefAuto{A\neq\varnothing \eqvdash \exists x\,(x \in A)}{5}}

  \proofstep{8}{a\in X}{\rA}
  \proofstep{2,8}{a\in A}{%
    \FormulaRefAuto{A\subseteq B,\, x\in A \vdash x\in B}{6,8}}
  \proofstep{2,8}{(X,a)\in \PowNE{A}\times A}{%
    \FormulaRefAuto{a\in A\dsep b\in B\vdash (a,b)\in A\times B}{2,9}}
  \proofstep{2,8}{(X,a)\in \ChoiceRel{A}}{%
    \rIE{\FormulaRefAuto{\ChoiceRel{A}\coloneqq
      \{(X,a)\in \PowNE{A}\times A\mid a\in X\}},
      \FormulaRefAuto{x \in A\dsep P(x)\vdash x \in \{x \in A \mid P(x)\}}{10,8}}}
  \proofstep{2,8}{\exists a\,\bigl((X,a)\in \ChoiceRel{A}\bigr)}{\rEI{11}}
  \proofstep{2}{\exists a\,\bigl((X,a)\in \ChoiceRel{A}\bigr)}{\rEE{7,8,12}}

  \proofstep{}{%
    X\in \PowNE{A}\rightarrow \exists a\,\bigl((X,a)\in \ChoiceRel{A}\bigr)%
  }{\rRI{2,13}}
  \proofstep{}{%
    \TotRel{\ChoiceRel{A},\PowNE{A},A}%
  }{\FormulaRefAuto{Totale Relation}{1,14}}
\end{tabproof}

\FormulaThmDeltaK[Elemente der Auswahlrelation]{%
  (X,a)\in \ChoiceRel{A}\vdash a\in X%
}{ChoiceRelElim}{
  \DeltaDecl{Mengen}{A \dsep X \dsep a}
  \DeltaDecl{Relation}{\ChoiceRel{A}}
}
\begin{tabproof}
  \proofstep{1}{(X,a)\in \ChoiceRel{A}}{\rA}
  \proofstep{1}{%
    (X,a)\in \{(X,a)\in \PowNE{A}\times A\mid a\in X\}%
  }{%
    \rIE{\FormulaRefAuto{\ChoiceRel{A}\coloneqq
      \{(X,a)\in \PowNE{A}\times A\mid a\in X\}},1}}
  \proofstep{1}{(X,a)\in \PowNE{A}\times A\land a\in X}{%
    \FormulaRefAuto{x \in \{u \in A \mid P(u)\} \eqvdash x \in A \land P(x)}{2}}
  \proofstep{1}{a\in X}{\rAEb{3}}
\end{tabproof}

\FormulaThmDeltaK[Auswahlrelation liefert Auswahlfunktion]{%
  \begin{aligned}[t]
    &\exists C\colon \PowNE{A}\to A\,
      \forall X\in \PowNE{A}\,
        \bigl((X,C(X))\in \ChoiceRel{A}\bigr)\vdash{}\\
    &\exists C\,\ChoicePow{C}{A}
  \end{aligned}%
}{ChoiceRelToChoicePow}{
  \DeltaDecl{Mengen}{A \dsep X}
  \DeltaDecl{Funktionen}{C}
  \DeltaDecl{Relation}{\ChoiceRel{A}}
}
\paragraph{Beweis.}
Sei \(C\) ein Zeuge der Voraussetzung. Dann gilt
\[
  C\colon\PowNE{A}\to A
  \quad\text{und}\quad
  \forall X\in\PowNE{A}\,
    \bigl((X,C(X))\in\ChoiceRel{A}\bigr).
\]
Für ein beliebiges \(X\in\PowNE{A}\) liefert
\FormulaRefAuto{ChoiceRelElim}[thm] daher \(C(X)\in X\). Folglich gilt
\(\forall X\in\PowNE{A}\,(C(X)\in X)\). Zusammen mit der Typisierung von
\(C\) folgt nach \FormulaRefAuto{ChoicePow}[def] die Aussage
\(\ChoicePow{C}{A}\), und die Existenzquantifizierung über \(C\) ergibt die
Behauptung.
\hfill\(\square\)

\FormulaThmDeltaK[Surjektionsform liefert Auswahlfunktion auf nichtleeren Teilmengen]{%
  \ACsur\vdash \exists C\,\ChoicePow{C}{A}%
}{ACsurChoicePow}{
  \DeltaDecl{Mengen}{A \dsep X}
  \DeltaDecl{Funktionen}{C \dsep s}
  \DeltaDecl{Relation}{\ChoiceRel{A}}
}
\begin{tabproof}
  \proofstep{1}{\ACsur}{\rA}
  \proofstep{1}{%
    \forall A\,\forall B\,\forall G\,
    \bigl(G\colon B\sur A\rightarrow
    \exists F\colon A\to B\,G\circ F=\Id_A\bigr)%
  }{\rIE{\FormulaRefAuto{ACsur},1}}
  \proofstep{}{\TotRel{\ChoiceRel{A},\PowNE{A},A}}{\FormulaRefAuto{ChoiceRelTotal}}
  \proofstep{}{\pi_1\restriction_{\ChoiceRel{A}}\colon \ChoiceRel{A}\sur \PowNE{A}}{%
    \FormulaRefAuto{\pi_1\restriction_R\colon R\sur A}{3}}
  \proofstep{1}{%
    \exists s\colon \PowNE{A}\to \ChoiceRel{A}\,
    \bigl(\pi_1\restriction_{\ChoiceRel{A}}\circ s=\Id_{\PowNE{A}}\bigr)%
  }{\rRE{\rUE{2},4}}
  \proofstep{1}{%
    \exists C\colon \PowNE{A}\to A\,
    \forall X\in \PowNE{A}\,\bigl((X,C(X))\in \ChoiceRel{A}\bigr)%
  }{%
    \FormulaRefAuto{\exists s\colon A\to R\,\bigl(\pi_1\restriction_R\circ s = \Id_A\bigr)
    \vdash \exists F\colon A\to B\,\forall x\in A\,\bigl((x,F(x))\in R\bigr)}{5}}
  \proofstep{1}{\exists C\,\ChoicePow{C}{A}}{\FormulaRefAuto{ChoiceRelToChoicePow}{6}}
\end{tabproof}

\FormulaThmDeltaK[Zermelos Konstruktionssatz]{%
  \ChoicePow{C}{A}\vdash
  \exists\,\preccurlyeq\;\WellOrd{A,\preccurlyeq}%
}{ZermeloConstruction}{
  \DeltaDecl{Mengen}{A}
  \DeltaDecl{Funktion}{C}
  \DeltaRow{Relationsoperatoren}{\preccurlyeq}
  \DeltaRow{Auswahlfunktion}
           {\ChoicePow{C}{A}}
           [\FormulaRefAuto{ChoicePow}]
  \DeltaRow{Konstruktionsresultat}
           {\exists\,\preccurlyeq\;\WellOrd{A,\preccurlyeq}}
}

\FormulaThmDeltaK[Zermelos Wohlordnungskonstruktion]{%
  \exists C\,\ChoicePow{C}{A}\vdash \WellOrdable{A}%
}{ZermeloWO}{
  \DeltaDecl{Mengen}{A}
  \DeltaDecl{Funktion}{C}
}
\begin{tabproof}
  \proofstep{1}{\exists C\,\ChoicePow{C}{A}}{\rA}
  \proofstep{2}{\ChoicePow{C}{A}}{\rA}
  \proofstep{2}{\exists\,\preccurlyeq\;\WellOrd{A,\preccurlyeq}}{%
    \FormulaRefAuto{ZermeloConstruction}{2}}
  \proofstep{2}{\WellOrdable{A}}{\rIE{\FormulaRefAuto{WellOrdable},3}}
  \proofstep{1}{\WellOrdable{A}}{\rEE{1,2,4}}
\end{tabproof}

\FormulaThmDeltaK[Auswahlaxiom impliziert Wohlordnungssatz]{%
  \ACsur\vdash \WOP%
}{ACsurToWOP}{
  \DeltaDecl{Mengen}{A}
}
\begin{tabproof}
  \proofstep{1}{\ACsur}{\rA}
  \proofstep{1}{\exists C\,\ChoicePow{C}{A}}{\FormulaRefAuto{ACsurChoicePow}{1}}
  \proofstep{1}{\WellOrdable{A}}{\FormulaRefAuto{ZermeloWO}{2}}
  \proofstep{1}{\forall A\,\WellOrdable{A}}{\rUI{3}}
  \proofstep{1}{\WOP}{\rIE{\FormulaRefAuto{WOP},4}}
\end{tabproof}

\FormulaDefDeltaK[Faserminimenge zu einer Wohlordnung]{%
  L_{G,\preccurlyeq}\coloneqq
  \{\,b\in B\mid
    \forall c\in B\,
      \bigl(G(c)=G(b)\rightarrow b\preccurlyeq c\bigr)\,\}%
}{FiberMinSet}{
  \DeltaDecl{Mengen}{A \dsep B\dsep b \dsep c}
  \DeltaDecl{Funktion}{G}
  \DeltaRow{Relationsoperatoren}{\preccurlyeq}
  \DeltaDecl{Mengen}{L_{G,\preccurlyeq}}
}

\FormulaThmDeltaK[Wohlgeordnete Fasern liefern eine Auswahlmenge]{%
  \WellOrd{B,\preccurlyeq}\dsep G\colon B\sur A
  \vdash
  \forall X\in \Fib_G[A]\,
    \exists! b\,(b\in X\land b\in L_{G,\preccurlyeq})%
}{WellOrdFiberChoiceSet}{
  \DeltaDecl{Mengen}{A \dsep B\dsep X \dsep b}
  \DeltaDecl{Funktion}{G}
  \DeltaRow{Relationsoperatoren}{\preccurlyeq}
  \DeltaRow{Faserminimenge}{L_{G,\preccurlyeq}}
    [\FormulaRefAuto{FiberMinSet}]
}

\FormulaThmDeltaK[Faserminima einer Wohlordnung liefern eine Rechtsinverse]{%
  \WellOrd{B,\preccurlyeq}\dsep G\colon B\sur A
  \vdash \exists F\colon A\to B\,G\circ F=\Id_A%
}{WellOrdFibersSection}{
  \DeltaDecl{Mengen}{A \dsep B}
  \DeltaDecl{Funktionen}{F \dsep G}
  \DeltaRow{Relationsoperatoren}{\preccurlyeq}
}
\begin{tabproof}
  \proofstep{1}{\WellOrd{B,\preccurlyeq}}{\rA}
  \proofstep{2}{G\colon B\sur A}{\rA}
  \proofstep{1,2}{%
    \forall X\in \Fib_G[A]\,
      \exists! b\,(b\in X\land b\in L_{G,\preccurlyeq})%
  }{\FormulaRefAuto{WellOrdFiberChoiceSet}{1,2}}
  \proofstep{1,2}{\exists F\colon A\to B\,G\circ F=\Id_A}{%
    \FormulaRefAuto{ChoiceSetRightInverse}{3}}
\end{tabproof}

\FormulaThmDeltaK[Wohlordnungssatz impliziert Auswahlaxiom]{%
  \WOP\vdash \ACsur%
}{WOPToACsur}{
  \DeltaDecl{Mengen}{A \dsep B}
  \DeltaDecl{Funktionen}{F \dsep G}
}
\begin{tabproof}
  \proofstep{1}{\WOP}{\rA}
  \proofstep{1}{\forall M\,\WellOrdable{M}}{\rIE{\FormulaRefAuto{WOP},1}}
  \proofstep{3}{G\colon B\sur A}{\rA}
  \proofstep{1}{\WellOrdable{B}}{\rUE{2}}
  \proofstep{1}{%
    \exists\,\preccurlyeq\;\WellOrd{B,\preccurlyeq}}{%
    \rIE{\FormulaRefAuto{WellOrdable},4}}

  \proofstep{6}{\WellOrd{B,\preccurlyeq}}{\rA}
  \proofstep{3,6}{\exists F\colon A\to B\,G\circ F=\Id_A}{%
    \FormulaRefAuto{WellOrdFibersSection}{6,3}}
  \proofstep{1,3}{\exists F\colon A\to B\,G\circ F=\Id_A}{\rEE{5,6,7}}

  \proofstep{1}{%
    G\colon B\sur A\rightarrow
    \exists F\colon A\to B\,G\circ F=\Id_A%
  }{\rRI{3,8}}
  \proofstep{1}{%
    \forall G\,\bigl(G\colon B\sur A\rightarrow
    \exists F\colon A\to B\,G\circ F=\Id_A\bigr)%
  }{\rUI{9}}
  \proofstep{1}{%
    \forall B\,\forall G\,\bigl(G\colon B\sur A\rightarrow
    \exists F\colon A\to B\,G\circ F=\Id_A\bigr)%
  }{\rUI{10}}
  \proofstep{1}{%
    \forall A\,\forall B\,\forall G\,\bigl(G\colon B\sur A\rightarrow
    \exists F\colon A\to B\,G\circ F=\Id_A\bigr)%
  }{\rUI{11}}
  \proofstep{1}{\ACsur}{\rIE{\FormulaRefAuto{ACsur},12}}
\end{tabproof}

\FormulaThmDeltaK[Äquivalenz von Auswahlaxiom und Wohlordnungssatz]{%
  \ACsur\eqvdash \WOP%
}{ACsurEquivWOP}{
  \DeltaDecl{Mengen}{A \dsep B}
  \DeltaDecl{Funktionen}{F \dsep G}
}
\begin{tabproof}
  \proofstep{1}{\ACsur}{\rA}
  \proofstep{1}{\WOP}{\FormulaRefAuto{ACsurToWOP}{1}}
  \proofstep{}{\ACsur\rightarrow \WOP}{\rRI{1,2}}

  \proofstep{4}{\WOP}{\rA}
  \proofstep{4}{\ACsur}{\FormulaRefAuto{WOPToACsur}{4}}
  \proofstep{}{\WOP\rightarrow \ACsur}{\rRI{4,5}}

  \proofstep{}{\ACsur\eqvdash \WOP}{\rLRI{3,6}}
\end{tabproof}

\chapter{Ordnungsbezogene Teilmengen}

\section{Hauptfilter}

\FormulaDefDeltaK[Hauptfilter eines Elements]{%
  \PrincipalFilter{A,\leq}{j}
  \coloneqq \{x\in A\mid j\leq x\}%
}{PrincipalFilter}{%
  \DeltaRow{Mengen}{A\dsep j\dsep x}
  \DeltaRow{Relationsoperatoren}{\leq}
  \DeltaRow{\textbf{Neues Symbol}}{}
  \DeltaPrem{Hauptfilter}{\PrincipalFilter{A,\leq}{j}}
}
\leavevmode\par\medskip

\FormulaThmDeltaK[Hauptfilter liegen im Ordnungsträger]{%
  \PrincipalFilter{A,\leq}{j}\subseteq A
}{PrincipalFilterSubsetCarrier}{%
  \DeltaRow{Mengen}{A\dsep j\dsep x}
  \DeltaRow{Relationsoperatoren}{\leq}
}
\begin{tabproof}
  \proofstep{1}{x\in\PrincipalFilter{A,\leq}{j}}{\rA}
  \proofstep{1}{x\in A\land j\leq x}{%
    \FormulaRefAuto{PrincipalFilter}[def]{1}}
  \proofstep{1}{x\in A}{\rAEa{2}}
  \proofstep{}{\PrincipalFilter{A,\leq}{j}\subseteq A}{%
    \FormulaRefAuto{A \subseteq B \coloneq
      \forall x\,(x\in A\rightarrow x\in B)}{\rUI{\rRI{1,3}}}}
\end{tabproof}


\chapter{Wegweiser zu Ordnungsbeispielen}

Die Bandnummern geben die Beweisabhängigkeit wieder; sie schließen
fachübergreifende Beispiele nicht aus. Band~12 untersucht die
Mitgliedschafts-, Teilmengen- und additive Ordnung der natürlichen Zahlen.
Band~13 erweitert diese Ordnung auf die ganzen Zahlen, Band~14 verwendet die
Inklusionsordnung auf Familien von Teilmengen, und Band~17 gewinnt aus jeder
Halbverbandsoperation eine partielle Ordnung. Die duale Ordnung und
Hauptfilter wurden bereits in diesem Band als allgemeine Konstruktionen
behandelt; die endlichen Teilordnungen folgen in Band~14.

\end{document}
